% main.tex -- DRAFT. Not for circulation until the TODO list at the end of this file is cleared.
% Build:  pdflatex main && bibtex main && pdflatex main && pdflatex main
\documentclass[11pt]{amsart}

\usepackage[T1]{fontenc}
\usepackage[utf8]{inputenc}
\usepackage{lmodern}
\usepackage{amsmath,amssymb,amsthm,mathtools}
\usepackage{booktabs}
\usepackage{array}
\usepackage{enumitem}
\usepackage{alltt}
\usepackage{xcolor}
\usepackage[margin=1.1in]{geometry}
\usepackage[colorlinks=true,linkcolor=blue!50!black,citecolor=green!40!black,urlcolor=blue!50!black]{hyperref}

% ------------------------------------------------------------------ macros
\newcommand{\Z}{\mathbb{Z}}
\newcommand{\Q}{\mathbb{Q}}
\newcommand{\F}{\mathbb{F}}
\newcommand{\C}{\mathbb{C}}
\newcommand{\Zl}{\mathbb{Z}_\ell}
\newcommand{\Ql}{\mathbb{Q}_\ell}
\newcommand{\Fl}{\mathbb{F}_\ell}
\newcommand{\Cl}{\mathbb{C}_\ell}
\newcommand{\SL}{\mathrm{SL}_2(\Z)}
\newcommand{\Sb}{\overline{S}}
\newcommand{\Mb}{\overline{M}}
\newcommand{\Ub}{\overline{U}}
\newcommand{\Db}{\overline{D}}
\newcommand{\Xb}{\overline{X}}
\newcommand{\Lb}{\overline{L}}
\newcommand{\cS}{\mathcal{S}}
\newcommand{\cT}{\mathcal{T}}
\newcommand{\cV}{\mathcal{V}}
\newcommand{\cE}{\mathcal{E}}
\newcommand{\fm}{\mathfrak{m}}
\newcommand{\leg}[2]{\left(\frac{#1}{#2}\right)}
\DeclareMathOperator{\rank}{rank}
\DeclareMathOperator{\ord}{ord}
\DeclareMathOperator{\Lie}{Lie}
\DeclareMathOperator{\charpoly}{charpoly}
\DeclareMathOperator{\Span}{span}
\newcommand{\w}{w}
\newcommand{\st}[1]{\textup{\textsf{[#1]}}}
\newcommand{\TODO}[1]{\textcolor{red}{\textbf{[TODO:} #1\textbf{]}}}
\newcommand{\lean}[1]{\texttt{#1}}

% ------------------------------------------------------------------ theorems
\theoremstyle{plain}
\newtheorem{theorem}{Theorem}[section]
\newtheorem{proposition}[theorem]{Proposition}
\newtheorem{lemma}[theorem]{Lemma}
\newtheorem{corollary}[theorem]{Corollary}
\newtheorem{conjecture}[theorem]{Conjecture}
\newtheorem{mainthm}{Theorem}
\renewcommand{\themainthm}{\Alph{mainthm}}
\theoremstyle{definition}
\newtheorem{definition}[theorem]{Definition}
\newtheorem{remark}[theorem]{Remark}
\newtheorem{computation}[theorem]{Computational result}
\newtheorem{question}[theorem]{Question}

\numberwithin{equation}{section}

\begin{document}

\title[Atkin--O'Brien congruences and the vanishing of $d_\ell$]{Atkin--O'Brien congruences for $\ell=17,19,23$, a formal proof, and the vanishing of the Boylan--Webb invariant $d_\ell$}

\author{Thomas DiFiore}
\email{tomdif@gmail.com}
% TODO: author list, affiliation, and a disclosure statement on the use of AI assistance in the
% development of the proofs and the formalization, as required by the target journal.

\date{Draft, \today}

\begin{abstract}
For a prime $\ell\ge 5$ let $P_\ell(b)=\sum_n p\big((\ell^b n+1)/24\big)q^{n/24}$. Folsom, Kent and Ono showed that the
$\Z/\ell^m$-modules spanned by the $P_\ell(b)$ stabilize; Boylan and Webb sharpened the stabilization index to
$b_\ell(m)\le 2m-1$, conditionally on the vanishing of an invariant $d_\ell$ that they computed to be $0$ for $\ell<1300$.
We report three results of different status.
(1) For $\ell\in\{13,17,19,23\}$, which are exactly the primes $\ell\ge 13$ with $\dim S_{\ell-1}(\SL)=1$, there is a
single $\kappa_\ell\in\Zl^\times$ such that
$p\big((\ell^{m+2}n+1)/24\big)\equiv\kappa_\ell\, p\big((\ell^m n+1)/24\big)\pmod{\ell^m}$ for all $m\ge1$. For $\ell=13$ this is
the Atkin--O'Brien theorem. For $\ell=17,19,23$ it improves $b_\ell(m)\le 2m-1$ to $b_\ell(m)\le m$; to our knowledge it is
new. The theorem is verified in Lean~4 with Mathlib, with no hypotheses and only the standard axioms. The formal proof
includes formal groups of Weierstrass curves, Lubin's quotient, the canonical subgroup, V\'elu's formulas, the Tate
curve, Deligne's congruence for the Hasse invariant, Igusa's simple-zeros theorem and Ramanujan's differential equations.
(2) We give a paper proof that $T_{\ell^2}$ is invertible modulo $\ell$ on $S_{(\ell-2)/2}(\SL,\nu_\eta^{-\ell})$ for every
prime $\ell\ge 13$. It uses Yifan Yang's Shimura-type isomorphism with level-$6$ newforms and a $\theta$-cycle argument
at weight $6$. It implies $d_\ell=0$ for all $\ell\ge 13$, which is the partition case of a conjecture of Belmont, Lee,
Musat and Trebat-Leder, and hence $b_\ell(m)\le 2m-1$ unconditionally.
(3) We explain Boylan and Webb's rank deficit at $\ell=607$ by a two-dimensional eigenspace of $T_{607^2}$ modulo $607$,
and we give a cyclicity criterion for attaining the Folsom--Kent--Ono rank bound. We also report computations
supporting $b_\ell(m)=2m-1$ for $29\le\ell\le 229$.
\end{abstract}

\maketitle
\setcounter{tocdepth}{1}
\tableofcontents

% =====================================================================
\section{Introduction}\label{sec:intro}
% =====================================================================

\subsection{Background}
Ramanujan~\cite{Ramanujan1919} discovered the congruences $p(5n+4)\equiv0\pmod5$, $p(7n+5)\equiv0\pmod7$ and
$p(11n+6)\equiv 0\pmod{11}$, and conjectured their extension to powers of $5$, $7$ and $11$. Watson~\cite{Watson1938}
proved the conjecture for powers of $5$ and a corrected form for powers of $7$, and Atkin~\cite{Atkin1967} proved it for
powers of $11$. For $\ell=13$ the behaviour is different. Atkin and O'Brien~\cite{AtkinOBrien1967} proved that there are
constants $K_a$ such that
\begin{equation}\label{eq:AO}
  p\!\left(\frac{13^{a+2}n+1}{24}\right)\equiv K_a\, p\!\left(\frac{13^{a}n+1}{24}\right)\pmod{13^{a}}
\end{equation}
(see also the restatement in Garvan--Shomanov~\cite[Thm.~1.1]{GarvanShomanov2025}, who also proved Atkin and O'Brien's
Hecke-type conjecture modulo powers of $13$). Atkin~\cite{Atkin1968} studied related multiplicative congruence properties
for other primes.

Folsom, Kent and Ono~\cite{FKO2012} placed these phenomena in a general $\ell$-adic framework. For every prime
$\ell\ge 5$, the $\Z/\ell^m$-modules generated by the generating functions of
$p\big((\ell^b n+1)/24\big)$ stabilize as $b$ grows within a parity class, and the stable modules have rank at most
$\lfloor(\ell+12)/24\rfloor$. Ramsey's appendix to~\cite{FKO2012} identifies the stable object with the ordinary part of
$U(\ell^2)$ on overconvergent half-integral-weight forms, which answers a question of Mazur. Boylan and
Webb~\cite{BoylanWebb2013} introduced an invariant $d_\ell$ of the mod-$\ell$ action and proved
$b_\ell(m)\le 2(d_\ell+1)m-1$~\cite[Thm.~5.1]{BoylanWebb2013}; in particular $d_\ell=0$ implies $b_\ell(m)\le 2m-1$~\cite[Thm.~1.2]{BoylanWebb2013}. They
verified $d_\ell=0$ by computer for $5\le\ell<1300$ and remarked that $2m-1$ is sharp in all computed cases with
$\ell\ge 29$. They also found that the Folsom--Kent--Ono rank bound is attained for every $\ell\le1297$ except
$\ell=607$, where the stable rank is one less~\cite[\S6]{BoylanWebb2013}. Belmont, Lee, Musat and Trebat-Leder~\cite{BLMT2014}
extended the framework to $r$-coloured partitions and conjectured that their invariant $d_\ell(r)$ vanishes for all
$\ell\ge r+5$~\cite[\S5]{BLMT2014}. Our second main result concerns the case $r=1$.

\subsection{Notation}\label{sec:notation}
We follow Boylan--Webb~\cite{BoylanWebb2013}. Throughout, $\ell\ge 5$ is prime, $q=e^{2\pi i z}$ and
$\eta(z)=q^{1/24}\prod_{n\ge1}(1-q^n)$. On series in $q^{1/24}$ let $U=U_\ell$ act by
$U\big(\sum a(n)q^{n/24}\big)=\sum a(\ell n)q^{n/24}$, and let $V(f)(z)=f(\ell z)$. Put
\[
  \Phi=\Phi_\ell=\frac{\eta(\ell^2z)}{\eta(z)},\qquad D(f)=U(\Phi f),\qquad X=D\circ U,
\]
where $f|X$ means $f|U|D$: apply $U$ first, then $D$. Define $L(0)=1$, and for $b\ge1$ let $L(b)=D L(b-1)$ when $b$ is odd and
$L(b)=U L(b-1)$ when $b$ is even. Then
\[
  L(b)=\begin{cases}\eta(\ell z)\,P_\ell(b)& b\text{ odd},\\ \eta(z)\,P_\ell(b) & b \text{ even},\end{cases}\qquad
  P_\ell(b)=\sum_{n\ge 0,\ 24\mid \ell^bn+1} p\!\left(\frac{\ell^b n+1}{24}\right)q^{n/24}.
\]
Since $\ell^2\equiv1\pmod{24}$, the series $P_\ell(b)$ and $P_\ell(b+2)$ are supported on the same exponents.

\begin{definition}[stabilization index]\label{def:b}
For $b\ge1$ and $m\ge1$ let $\Lambda(b,m)$ be the $\Z/\ell^m$-span of the reductions modulo $\ell^m$ of
$\{L(b'):b'\ge b,\ b'\equiv b\ (2)\}$. Let $b_\ell(m)$ be the least $b\ge 1$ such that $\Lambda(b',m)=\Lambda(b'+2,m)$
for all $b'\ge b$. Let $b^{\mathrm{odd}}_\ell(m)$ be the least odd $b$ such that $X:\Lambda(b,m)\to\Lambda(b+2,m)$ is an
isomorphism. By~\cite[Remark~1]{BoylanWebb2013}, this is the threshold in Boylan and Webb's Theorem~1.4.
The rank $r_\ell(m)$ is the rank of the stable odd module $\Omega^{\mathrm{odd}}(m)$, and
$R_\ell:=\lfloor(\ell+12)/24\rfloor$ is the Folsom--Kent--Ono bound $r_\ell(m)\le R_\ell$~\cite[(1.13)]{BoylanWebb2013}.
\end{definition}

\begin{remark}[conventions]\label{rem:conventions}
The parity-free index $b_\ell(m)$ is our reconstruction. It reproduces Boylan and Webb's values $b_{29}(2)=3$ and
$b_{29}(3)=5$, and it agrees with the proportionality formulation whenever the stable module is free of rank one.
\TODO{Check Definition~\ref{def:b} line by line against \cite[\S1, Thm.~1.4 and Remark~1]{BoylanWebb2013}. In particular,
check whether Boylan--Webb's $b_\ell(m)$ is defined using both parity chains or only the odd one.}
\end{remark}

Let $d=d_\ell=\dim S_{\ell-1}(\SL)$. This $d$ is the dimension, not the Boylan--Webb invariant; the invariant is always
written $d_\ell$ with the subscript and is defined in~\eqref{eq:dl} below. Thus $d=1$ exactly for
$\ell\in\{13,17,19,23\}$, and $d=2$ for $\ell\in\{29,31\}$.
\TODO{$d_\ell$ is used for both quantities in the literature; rename the dimension, e.g.\ to $\delta_\ell$, throughout.}

\subsection{Main results}

\begin{mainthm}[Atkin--O'Brien congruences for $\dim S_{\ell-1}=1$; formally verified]\label{thm:A}
Let $\ell\in\{13,17,19,23\}$ and $\kappa^{(0)}_\ell=6,6,10,5$ respectively. There is $\kappa_\ell\in\Zl$ with
$\kappa_\ell\equiv\kappa_\ell^{(0)}\pmod\ell$ such that for every $m\ge1$ and every $n\in\Z_{\ge0}$ with
$24\mid\ell^m n+1$,
\begin{equation}\label{eq:main}
  \ell^m \ \Big|\ p\!\left(\frac{\ell^{m+2}n+1}{24}\right)-\kappa_\ell\; p\!\left(\frac{\ell^{m}n+1}{24}\right)\quad\text{in }\Zl .
\end{equation}
Equivalently, $P_\ell(b+2)\equiv\kappa_\ell P_\ell(b)\pmod{\ell^m}$ for all $b\ge m\ge1$.
\end{mainthm}

\noindent\emph{Status.} \st{Verified in Lean}. These are the theorems \lean{OrdQuot.atkin\_obrien\_13},
\lean{\_17}, \lean{\_19} and \lean{\_23} in \lean{RamanujanTau/OrdQuot/Atkin.lean}. They have no hypotheses, and their
dependency closure uses only the axioms \lean{propext}, \lean{Classical.choice} and \lean{Quot.sound}
(\S\ref{sec:lean}). The mathematics and the formal statement still need review by human specialists.

\medskip
The constant $\kappa_\ell$ does not depend on $m$: it is a single $\ell$-adic number, the unit eigenvalue of $X$ on the
stable line. For $\ell=13$, the reduction of $\kappa_{13}$ modulo $13^2$ is Atkin--O'Brien's $K_2=45$
(Table~\ref{tab:kappa}).

\begin{corollary}[translation to stabilization indices]\label{cor:bm}
For $\ell\in\{13,17,19,23\}$ and $m\ge1$, the stable modules are free of rank one, and $b_\ell(m)\le m$. The odd and even
chains are separately stable from the least $b\ge m$ of the respective parity. In particular
$b^{\mathrm{odd}}_\ell(m)\le m$ for odd $m$ and $b^{\mathrm{odd}}_\ell(m)\le m+1$ for even $m$.
\end{corollary}
\noindent\emph{Status.} \st{Proved}: a short deduction from Theorem~\ref{thm:A} (\S\ref{sec:translation}). The deduction
itself is not formalized. Equality $b_\ell(m)=m$ is \st{computed} for $m\le 16$ in a validated truncated model, and
directly from $p(n)$ data for $m\le 4$ ($\ell=17$) and $m\le 3$ ($\ell=13,19,23$) (\S\ref{sec:numerics}).

\medskip
For $\ell\ge 13$ write $j\equiv\ell\pmod{24}$ with $j\in\{1,5,7,11,13,17,19,23\}$, $t=\lfloor\ell/24\rfloor$ and
$\lambda=(\ell-3)/2$, and let $\nu_\eta$ be the multiplier system of $\eta$. Put
\[
  \cS_\ell:=S_{\lambda+\frac12}\big(\SL,\nu_\eta^{-\ell}\big).
\]
This space has dimension $R_\ell$ (Lemma~\ref{lem:dims}) and carries the Hecke operator $T_{\ell^2}$, which is integral.

\begin{mainthm}[Theorem A of the project docs; paper proof]\label{thm:B}
For every prime $\ell\ge13$, $\det\big(T_{\ell^2}\,|\,\cS_\ell\big)$ is an $\ell$-adic unit.
\end{mainthm}

\begin{corollary}\label{cor:dl}
For every prime $\ell\ge13$:
\begin{enumerate}[label=(\roman*)]
  \item $\Db(\Mb_{\ell-1})=\Db(\Sb_{\ell-1})=\Delta^t S_{k'}(\Fl)$, where $k'=\ell-1-12t$, and
        $\rank\Db=\rank\Xb=R_\ell=\lceil d/2\rceil$;
  \item $\Sb_{\ell-1}=\Delta^tS_{k'}(\Fl)\oplus\ker\Xb$, so $\Xb$ is semisimple at $0$;
  \item the Boylan--Webb invariant vanishes: $d_\ell=0$;
  \item $b_\ell(m)\le 2m-1$ for all $m\ge1$, by~\cite[Thm.~1.2]{BoylanWebb2013};
  \item $X$ has exactly $R_\ell$ unit roots on the odd overconvergent space.
\end{enumerate}
\end{corollary}
\noindent\emph{Status of Theorem~\ref{thm:B} and Corollary~\ref{cor:dl}.} \st{Paper proof}, citing
Yang~\cite[Thm.~1]{Yang2014}, the Deligne--Serre lifting lemma~\cite{DeligneSerre1974}, the theory of $\theta$-cycles
at level $N$~\cite{Jochnowitz1982,Katz1977,Edixhoven1992}, and the structure of the three-dimensional space
$S_6(\Gamma_0(6))$. None of this is formalized. Yang's theorem was checked against the arXiv version of his paper
(Remark~\ref{rem:yangcheck}). Item~(iii) is the case $r=1$ of the conjecture of~\cite[\S5]{BLMT2014}; to our knowledge
it has not appeared in the literature. Items (i)--(iii) were computed independently for all primes
$29\le\ell\le 229$ and for $\ell=607$ and $\ell=859$ (\S\ref{sec:numerics}). At $\ell=859$ the level-$6$ newform
$6.6.a.a$ is non-ordinary, so a weaker version of the sign argument would fail there.

\medskip
Our third group of results concerns attainment of the Folsom--Kent--Ono rank bound and sharpness of $2m-1$
(\S\S\ref{sec:rank}--\ref{sec:numerics}). Assuming Theorem~\ref{thm:B}:
\begin{itemize}
\item \st{Proved}. If $\Lb(1)=\eta(\ell z)U_\ell(1/\eta)\bmod\ell$ is a cyclic vector for $T_{\ell^2}$ on
$\cS_\ell\otimes\Fl$, then the stable modules are free of rank $R_\ell$ for every $m$ (Proposition~\ref{prop:nakayama}).
The mod-$\ell$ rank $r_\ell(1)$ equals the Krylov rank of $\Lb(1)$. The converse for $m\ge2$ is open
(\S\ref{sec:rank}, gap G6).
\item \st{Computed}. At $\ell=607$, $T_{607^2}$ has a two-dimensional eigenspace modulo $607$ (eigenvalue $-367$), so no
vector is cyclic and $r_{607}(1)=R_{607}-1=24$. This explains Boylan--Webb's exceptional prime by an eigenvalue
coincidence. It does not contradict $d_{607}=0$, which also holds.
\item \st{Computed}, with a \st{conditional} extension. $b_\ell(m)=2m-1$ for every prime $29\le\ell\le229$ and every
$m$ up to $16$, $9$, $8$ or $7$ depending on $\ell$ (Table~\ref{tab:sweep}). For $29\le\ell\le113$ the computations
also verify a sufficient condition for $b_\ell(m)=2m-1$ for \emph{all}~$m$. The deduction depends on a truncation
argument that has not yet been written out (\S\ref{sec:theta}).
\end{itemize}

\subsection{Status at a glance}
\begin{center}
\small
\begin{tabular}{@{}p{0.47\textwidth}p{0.18\textwidth}p{0.26\textwidth}@{}}
\toprule
Statement & Status & Where\\
\midrule
Thm~\ref{thm:A}: $b_\ell(m)\le m$, $\ell\in\{13,17,19,23\}$ & Lean, unconditional & \S\S\ref{sec:proofA}--\ref{sec:lean}\\
Equality $b_\ell(m)=m$ for those $\ell$ & computed ($m\le16$) & \S\ref{sec:numerics}\\
Thm~\ref{thm:B}: $T_{\ell^2}$ invertible mod $\ell$, $\ell\ge13$ & paper proof & \S\ref{sec:thmB}\\
Cor.~\ref{cor:dl}: $d_\ell=0$, $b_\ell(m)\le2m-1$ for all $\ell$ & paper proof & \S\ref{sec:thmB}\\
Cyclic $\Lb(1)$ $\Rightarrow$ $r_\ell(m)=R_\ell$ for all $m$ & paper proof (given Thm~\ref{thm:B}) & \S\ref{sec:rank}\\
Non-cyclic $\Rightarrow$ deficit at every $m$ & open (G6) & \S\ref{sec:rank}\\
$\ell=607$: 2-dim.\ eigenspace of $T_{607^2}$ & computed & \S\ref{sec:607}\\
$b_\ell(m)=2m-1$, $29\le\ell\le229$, $m\le 7$--$16$ & computed & \S\ref{sec:numerics}\\
$b_\ell(m)=2m-1$ for all $m$, $29\le\ell\le113$ & computed + unwritten truncation step & \S\ref{sec:theta}\\
$b_\ell(m)=2m-1$ for all $\ell\ge 29$ & conjectural & \S\ref{sec:open}\\
\bottomrule
\end{tabular}
\end{center}

\subsection{Relation to the literature}
Theorem~\ref{thm:A} for $\ell=13$ is the Atkin--O'Brien theorem~\cite{AtkinOBrien1967}. Our proof is different: it
uses overconvergent forms and Katz expansions in the spirit of~\cite{Katz1973,Coleman1996,Buzzard2003,Vonk2015}, not
Atkin--O'Brien's explicit modular equations. To our knowledge, the cases $\ell=17,19,23$ have not been stated or proved
before, although~\cite{BoylanWebb2013} gives no table of $b_\ell(m)$ for $13\le\ell\le23$ and Atkin~\cite{Atkin1968}
records related numerical observations.
\TODO{Confirm with specialists and against~\cite{Atkin1968} that the cases $\ell=17,19,23$ of Theorem~\ref{thm:A} are
new. The literature check (docs/LIT\_CHECK.md) found no statement in~\cite{BoylanWebb2013,FKO2012,BLMT2014,Webb2011,ElGuindyGhazy2024}.}
The vanishing $d_\ell=0$ for all $\ell\ge13$ (Corollary~\ref{cor:dl}) was conjectured, for $r$-coloured partitions,
in~\cite{BLMT2014}, which calls the question of when $d_\ell(r)=0$ an open problem. Our proof covers only $r=1$.
Lemma~\ref{lem:imageD} is the Folsom--Kent--Ono bound $r_\ell(m)\le R_\ell$ in a one-line form. It is essentially
Boylan--Webb's Theorem~4.3 and Lemma~4.4, which state it as the vanishing order
$\ord_\infty\ge\lceil\ell/24\rceil$~\cite{BoylanWebb2013}.

\subsection*{Acknowledgements and provenance}
Parts of the arguments in \S\ref{sec:I2} originate in an outside proposal; we corrected it in three places
(Remark~\ref{rem:provenance}). Two independent external reviews checked the Yang citation and recomputed the case
$\ell=859$. The Lean proof of Ramanujan's differential equations adapts code from Sphere-Packing-Lean~\cite{SpherePackingLean}
(Apache~2.0). \TODO{Names of the external reviewers, if they agree to be named; funding; AI-assistance disclosure.}

% =====================================================================
\section{The overconvergent set-up}\label{sec:setup}
% =====================================================================

In this section and the next, $\ell\ge13$ is prime unless stated otherwise. Let $E=E_{\ell-1}\in M_{\ell-1}(\Zl)$ be the
normalized Eisenstein series, so that $E\equiv1\pmod\ell$ (von Staudt--Clausen). Write $v$ for the valuation on $\Cl$
with $v(\ell)=1$.

\subsection{Katz expansions}
For $i\ge0$ choose a saturated $\Zl$-submodule $B_i\subset M_{i(\ell-1)}(\Zl)$ complementary to
$E\cdot M_{(i-1)(\ell-1)}(\Zl)$. Thus $B_0=\Zl$, and since $M_{\ell-1}=\Zl E\oplus S_{\ell-1}$ we may take
$B_1=S_{\ell-1}(\Zl)$. Any weight-$0$ overconvergent form $f$ has a unique \emph{Katz expansion}
\[
  f=\sum_{i\ge0}\frac{b_i}{E^i},\qquad b_i\in B_i\otimes\Ql,
\]
where $b_i\to0$ suitably~\cite[\S2.6]{Katz1973}. Fix bases of the $B_i$ and let $v_i(f)$ be the minimum valuation of
the coordinates of $b_i$. For $r\in\Cl$ with $0<v(r)<1/(\ell+1)$ put
\[
  B_E(r)=\{f:\ v_i(f)\ge i\,v(r)\ \text{for all } i\}.
\]
This is a ring, and it consists of the forms that are integral on the region $\{v(E_{\ell-1})\le v(r)\}$. The operators
$U$ and $V$ act as usual; on the Tate curve, $V$ is $q\mapsto q^\ell$.

\begin{proposition}[Step (K): Katz coefficients versus the sup norm]\label{prop:K}
Let $0<v(r)<1$. A Katz series $f=\sum_i b_i/E^i$ satisfies $v_i(f)\ge i\,v(r)$ for all $i$ if and only if
$|f(x)|\le1$ at every point $x=(E',\omega)$ with $v(E_{\ell-1}(x))\le v(r)$.
\end{proposition}
\begin{proof}[Sketch]
Katz's integral model of the region is $A[Y]/(E\cdot Y-r)$~\cite[\S\S2.6--2.8]{Katz1973}. Its special fibre $\bar E\cdot Y=0$ is
reduced because the Hasse invariant has simple zeros (Igusa~\cite{Igusa1958}). Integrality on the model is therefore
equivalent to the pointwise bound.
\end{proof}

In the formalization, Proposition~\ref{prop:K} is proved for Katz series whose pieces are isobaric polynomials in
$E_4,E_6$ (\lean{Igusa.katz\_iff\_hasse'}). The test points are good-reduction Weierstrass curves over $\Cl$, evaluated at
$(c_4,-c_6)$. The proof goes through the existence of supersingular points, Igusa's theorem (derived from the Serre
derivative and Ramanujan's differential equations), and a monomial complement $B_i$.

\subsection{The four inputs}\label{sec:inputs}
Let $s=(\ell+1)/2$, $\delta=(\ell^2-1)/24$, and define
\[
  W=\frac{\eta(\ell^2z)}{\eta(z)^{\ell^2}},\qquad Q=E\,\frac{\eta(\ell z)^2}{\eta(z)^{2\ell}},\qquad G=\frac{E}{V(E)},
\]
so that $\Phi=\Delta^\delta W$. By Newman's criterion~\cite{Newman1959}, $W$ is a form of weight $-(\ell^2-1)/2$ on
$\Gamma_0(\ell^2)$: here $\sum_\delta \delta r_\delta=0$, $\sum_\delta(\ell^2/\delta)r_\delta=1-\ell^4\equiv0\pmod{24}$ and
$\prod_\delta\delta^{r_\delta}=\ell^2$. Hence $E^sW$ is a genuine weight-$0$ function on $X_0(\ell^2)$.

\begin{itemize}
\item[(I1)] \emph{$U$-contraction.} For $0<v(r)<1/(\ell+1)$, $\ell\,U(B_E(r))\subseteq B_E(r^\ell)$, and
  $G\in1+\fm B_E(r)$. (Katz~\cite{Katz1973}; Wan~\cite{Wan1998}; Vonk~\cite[Thm.~B]{Vonk2015}.)
\item[(I2)] \emph{Sharp radius for $Q$.} $Q\in 1+\fm B_E(r^\ell)$; equivalently $v_i(Q)\ge\lceil\ell i/(\ell+1)\rceil$ and $Q\equiv1\pmod\ell$.
\item[(I3)] \emph{Frobenius.} $V$ maps $B_E(r^\ell)$ into $B_E(r)$.
\item[(I4)] \emph{Base case.} For $\ell\in\{13,17,19,23\}$: $L(b+2)\equiv\kappa^{(0)}_\ell L(b)\pmod\ell$ for all $b\ge1$,
  and each $L(b)$, $b\ge1$, is cuspidal.
\end{itemize}

From (I1), $b/E^j\in r^{-j}B_E(r)$ for $b\in B_j$. Letting $v(r)\uparrow1/(\ell+1)$, we obtain
\begin{equation}\tag{UB}\label{eq:UB}
  v_i\big(U(b/E^j)\big)\ \ge\ \Big\lceil\frac{\ell i-j}{\ell+1}\Big\rceil-1\qquad(b\in B_j).
\end{equation}
This bound agrees exactly with direct numerical computation for $i,j\le14$ at $\ell=17$.

% ---------------------------------------------------------------------
\subsection{Proof of (I3)}\label{sec:I3}
Let $x=(E',\omega)$ satisfy $v:=v(E_{\ell-1}(x))\le v(r)<1/(\ell+1)$, and let $H$ be its canonical subgroup, with
$\pi:E'\to E'/H$ and dual $\check\pi$. Then $Vf(x)=f(E'/H,\check\pi^*\omega)$. By Katz~\cite[3.10.7]{Katz1973}
(see Buzzard~\cite[Thm.~3.3(ii)]{Buzzard2003}), $v(E_{\ell-1}(E'/H))=\ell v\le v(r^\ell)$. Hence
$|Vf(x)|\le1$ whenever $f\in B_E(r^\ell)$; since $f$ has weight $0$, the choice of differential does not matter.
Proposition~\ref{prop:K} gives $Vf\in B_E(r)$. \qed

% ---------------------------------------------------------------------
\subsection{Proof of (I2)}\label{sec:I2}
Keep the notation of \S\ref{sec:I3}, but now allow $0<v<\ell/(\ell+1)$.
\begin{enumerate}[label=(\arabic*)]
\item \emph{Newton polygon.} The $[\ell]$-series of the formal group has the shape
  $[\ell](T)=\ell T+\cdots+aT^\ell+\cdots+uT^{\ell^2}$ with $v(a)=v$ and $v(u)=0$. If $v<\ell/(\ell+1)$, the polygon
  breaks: the $\ell-1$ nonzero points of $H$ have valuation $(1-v)/(\ell-1)$, and the other $\ell^2-\ell$ nonzero torsion
  points have valuation $v/(\ell^2-\ell)$ (proof of~\cite[3.10.7]{Katz1973}).
\item \emph{Derivatives.} Lubin's parameter $Y=\prod_{P\in H}T(X\oplus P)$~\cite{Lubin1967} gives a N\'eron generator
  $\omega_0'$ on $E'/H$ with $\pi^*\omega_0'=\mu\omega$, where $v(\mu)=\sum_{P\ne0}v(T(P))=1-v$. Since
  $\pi^*\check\pi^*=\ell$, we get $\check\pi^*\omega=\lambda\omega'_0$ with $v(\lambda)=v$.
  As a cross-check, $\ker\check\pi=E'[\ell]/H$. For $v<1/(\ell+1)$ it is non-canonical in $E'/H$, which gives
  $v(\lambda)=(\ell-1)\ell v/(\ell^2-\ell)=v$. For $1/(\ell+1)<v<\ell/(\ell+1)$ it is canonical, which gives
  $v(\lambda)=1-(1-v)=v$.
\item \emph{Pointwise value.} $\dfrac{V\Delta}{\Delta^\ell}(x)=\lambda^{-12}\dfrac{\Delta(E'/H,\omega_0')}{\Delta(E',\omega)^\ell}$,
  and both $\Delta$-values are units, so $v\big(V\Delta/\Delta^\ell(x)\big)=-12v$.
\item \emph{Conclusion.} The $q$-series identity $Q^{12}=E^{12}\,V\Delta/\Delta^\ell$ gives $|Q(x)|^{12}=|\ell|^{12v-12v}=1$
  on the whole region $\{v(E)\le\rho\}$, for every $\rho<\ell/(\ell+1)$. The Hasse lift $A_\ell=\eta^{2\ell}/\eta(\ell z)^2$ is
  analytic and nonvanishing there, so $Q=E/A_\ell$ is analytic. Proposition~\ref{prop:K} at every such radius gives
  $v_i(Q)\ge i\rho$ for all $\rho<\ell/(\ell+1)$, i.e.\ $v_i(Q)\ge\lceil\ell i/(\ell+1)\rceil$.
\item \emph{Congruence.} $Q\equiv1\pmod\ell$ as a $q$-series, because $A_\ell\equiv1$ by Frobenius and $E\equiv1$. Mod-$\ell$
  $q$-expansions are injective on $\bigoplus_i\bar B_i$, so $Q\in1+\fm B_E(r^\ell)$ for $v(r)<1/(\ell+1)$. \qed
\end{enumerate}

\begin{remark}[provenance]\label{rem:provenance}
The outline of this argument comes from an outside proposal. We corrected three points: the sign in step~(3) (the
proposal had $+12v$); the relation between $Q^{12}$ and the proposal's $\hat Q$; and a spurious factor
$\Delta^{\ell^2-1}$ in the proposal's $\hat G$. The correct twelfth-power identity is
$(E^sW)^{24}=\hat Q^\ell V(\hat Q)(E/VE)^{12}$ with $\hat Q=Q^{12}=E^{12}V(\Delta)/\Delta^\ell$.
The values $v_i(Q)$ computed numerically agree with $\lceil\ell i/(\ell+1)\rceil$ exactly. Katz~3.10.7 and Buzzard~3.3
were checked against Buzzard's text. Katz~2.6.2, 3.11.2, 3.11.3 and Appendix~2, and Coleman~\cite{Coleman1996}, p.~221,
were not. \TODO{Check these four citations against the sources.}
\end{remark}

% ---------------------------------------------------------------------
\subsection{The $D$-bound}
\begin{lemma}\label{lem:eta}
$(E^sW)^2=Q^\ell\cdot V(Q)\cdot G$.
\end{lemma}
\begin{proof}
$Q^\ell V(Q)G=E^\ell\frac{\eta(\ell z)^{2\ell}}{\eta(z)^{2\ell^2}}\cdot V(E)\frac{\eta(\ell^2z)^2}{\eta(\ell z)^{2\ell}}\cdot\frac{E}{V(E)}
=E^{\ell+1}\frac{\eta(\ell^2 z)^2}{\eta(z)^{2\ell^2}}=(E^sW)^2.$
\end{proof}

\begin{proposition}\label{prop:DB}
$E^sW\in B_E(r)$ for $0<v(r)<1/(\ell+1)$. Consequently, for $b\in B_j$,
\begin{equation}\tag{DB}\label{eq:DB}
  v_i\big(D(b/E^j)\big)\ \ge\ \Big\lceil\frac{\ell i-j-s}{\ell+1}\Big\rceil-1 .
\end{equation}
\end{proposition}
\begin{proof}
By (I1)--(I3), $Q^\ell\,V(Q)\,G\in1+\fm B_E(r)$. The binomial coefficients $\binom{1/2}{n}$ are $\ell$-integral, and $q$-expansions
determine Katz expansions, so the square root $E^sW$ lies in $B_E(r)$. Since $12\delta=s(\ell-1)$, the quotient
$\Delta^\delta/E^s$ lies in $r^{-s}B_E(r)$, and
$\Phi\cdot b/E^j=(\Delta^\delta/E^s)(E^sW)(b/E^j)\in r^{-(j+s)}B_E(r)$. Applying (I1) and letting
$v(r)\uparrow1/(\ell+1)$ gives~\eqref{eq:DB}.
\end{proof}
Numerically, $v_i(E^sW)\ge\lceil i/(\ell+1)\rceil$ holds for $i\le64$, and~\eqref{eq:DB} matches the computed values for
$i,j\le12$ at $\ell=17$.

% =====================================================================
\section{Proof of Theorem~\ref{thm:A}}\label{sec:proofA}
% =====================================================================

\subsection{Envelopes}
Define
\[
  \alpha(i)=\Big\lceil\frac{\ell i-2}{\ell+1}\Big\rceil,\qquad \beta(i)=\Big\lceil\frac{\ell i-2-s}{\ell+1}\Big\rceil,
\]
and, for an envelope $\gamma\in\{\alpha,\beta\}$ and $c\ge0$,
\[
  A_c(\gamma)=\{f\text{ cuspidal}:\ v_i(f)\ge\max(0,\gamma(i)-c)\text{ for all }i\}.
\]
Note that $\alpha(0)=\beta(0)=0$, that $\alpha(1)=\beta(1)=1$, and that $\alpha(i),\beta(i)\ge2$ for $i\ge2$ when $\ell\ge 13$.

\begin{lemma}[envelope lemma]\label{lem:env}
For every prime $\ell\ge13$, the bounds~\eqref{eq:UB} and~\eqref{eq:DB} imply:
\begin{enumerate}[label=(\arabic*)]
\item $U(A_2(\beta))\subseteq A_1(\alpha)$ and $D(A_2(\alpha))\subseteq A_1(\beta)$;
\item $U(A_1(\beta))\subseteq A_1(\alpha)$ and $D(A_1(\alpha))\subseteq A_1(\beta)$.
\end{enumerate}
\end{lemma}
\begin{proof}
By the ultrametric inequality, $v_i(Uf)\ge\min_j\big(v_j(f)+u(i,j)\big)$, where $u(i,j)$ is the right-hand side
of~\eqref{eq:UB}; the same holds for $D$ with~\eqref{eq:DB}. Split the minimum. For $j\le2$ the shift is exact. For
$j\ge3$ the slack $(\ell-1)j/(\ell+1)-c$ exceeds the loss. The tightest case is the $U$-step at $j=3$, which needs
$\big(3(\ell-1)-(\ell+1)/2\big)/(\ell+1)>2$, i.e.\ $\ell>11$.
\end{proof}
This lemma is formalized for every prime $\ell\ge13$ and all Katz indices: \lean{AtkinOBrien.Envelope.U\_step},
\lean{D\_step} and \lean{envelopes\_ge\_two}, assembled in \lean{AtkinOBrien.Bridge}.

\subsection{The base case and the role of $\dim S_{\ell-1}=1$}\label{sec:base}
\begin{lemma}\label{lem:base}
Let $\ell\in\{13,17,19,23\}$. For every $b\ge1$, $\Lb(b)$ is a nonzero element of the line $\Sb_{\ell-1}=S_{\ell-1}(\Fl)$, and
$\Lb(b+2)=\kappa^{(0)}_\ell\Lb(b)$ with $\kappa^{(0)}_\ell=6,6,10,5$.
\end{lemma}
\begin{proof}
By Lemma~\ref{lem:Uinv} below, $\Ub$ is invertible on $\Sb_{\ell-1}$. By Lemma~\ref{lem:imageD}, $\Db(\Mb_{\ell-1})$ lies in a
space of dimension $R_\ell=1=d$ inside $\Sb_{\ell-1}$. So $\Ub$ and $\Db$ act on the line $\Sb_{\ell-1}$ by scalars $u,\delta'$, and
$\Lb(1)=\Db(1)\in\Sb_{\ell-1}$. Then both $\Xb=\Db\Ub$ and $\Ub\Db$ act by $u\delta'$, and $\Lb(b+2)=u\delta'\,\Lb(b)$ for $b\ge1$.
Nonvanishing of $\Lb(1)$ and of $u\delta'$, and the values of $\kappa^{(0)}_\ell=u\delta'$, are finite computations. They are
certified in Lean by kernel-checked certificates (\lean{decide +kernel}) together with a formalized Sturm
bound~\cite{Sturm1987} (\lean{BoylanWebb.two\_step\_mod\_ell}, \lean{hnz\_\ensuremath{\ell}}, \lean{CertX.rank\_one}).
\end{proof}

\begin{remark}\label{rem:kappa-yang}
\st{Computed}. Under the dictionary of \S\ref{sec:dict}, $\kappa^{(0)}_\ell$ is the eigenvalue of $T_{\ell^2}$ on the line
$\cS_\ell\otimes\Fl$. Yang's isomorphism (Theorem~\ref{thm:yang}) then predicts
$\kappa^{(0)}_\ell\equiv\leg{12}{\ell}a_\ell(F_\ell)\pmod\ell$, where $F_\ell$ is the unique newform in
$S^{\mathrm{new}}_{\ell-3}(\Gamma_0(6))$ with Atkin--Lehner signs $(w_2,w_3)=(-\leg{2}{\ell},-\leg{3}{\ell})$. A PARI/GP
computation~\cite{PARI} confirms this: the class has dimension one in each case, and
$\leg{12}{\ell}a_\ell(F_\ell)\bmod\ell=6,6,10,5$ for $\ell=13,17,19,23$.
\end{remark}

\subsection{The induction, with constant absorption}
The following abstract statement is \lean{BoylanWebb.induction\_seq}. Let $\varpi$ be an element of a commutative ring $\mathcal{O}$,
and let $M$ be a $\varpi$-torsion-free $\mathcal O$-module. Suppose we are given linear maps $T_b:M\to M$ with $T_{b+2}=T_b$,
elements $L(b)\in M$ with $L(b+1)=T_bL(b)$, and submodules $A_1(b)$, $A_2(b)$, $2$-periodic in $b$, such that for all $b\ge1$:
\begin{enumerate}[label=(H\arabic*)]
\item\label{H:L} $L(b)\in A_1(b)$;
\item\label{H:T} $T_b(A_2(b))\subseteq A_1(b+1)$;
\item\label{H:lead} every $x\in A_1(b)$ can be written $x=c\,L(b)+\varpi y$ with $c\in\mathcal O$ and $y\in M$;
\item\label{H:div} if $x\in A_1(b)$ and $x=\varpi y$, then $y\in A_2(b)$;
\item\label{H:base} $L(\beta+2)-\kappa L(\beta)\in\varpi M$ for all $\beta\ge1$.
\end{enumerate}

\begin{theorem}[constant absorption]\label{thm:induction}
Under \ref{H:L}--\ref{H:base} there are $K_0=\kappa,K_1,K_2,\dots\in\mathcal O$ with
$K_{j+1}-K_j\in\varpi^{j+1}\mathcal O$, such that $L(\beta+2)-K_jL(\beta)\in\varpi^{j+1}M$ for all $\beta\ge j+1$.
If $\mathcal O=\Zl$ and $\varpi=\ell$, then $\kappa_\infty=\lim K_j$ satisfies $\kappa_\infty\equiv\kappa\pmod\ell$ and
$L(b+2)\equiv\kappa_\infty L(b)\pmod{\ell^b}$ for every $b\ge1$.
\end{theorem}
\begin{proof}
Write $\mathbf P(b,k)$ for the statement: for all $\beta\ge b$, $L(\beta+2)-kL(\beta)=\varpi^b g_\beta$ with $g_\beta\in A_2(\beta)$.
\ref{H:base} and \ref{H:div} give $\mathbf P(1,\kappa)$; note that $L(\beta+2)-kL(\beta)\in A_1(\beta)$ by \ref{H:L} and
periodicity. Assume $\mathbf P(b,k)$. Applying $T_\beta$ gives
$L(\beta+3)-kL(\beta+1)=\varpi^bT_\beta g_\beta$. Since $M$ is $\varpi$-torsion free, this forces $g_{\beta+1}=T_\beta g_\beta$.
By \ref{H:T} and \ref{H:lead}, $g_{b+1}=d\,L(b+1)+\varpi y_{b+1}$ for some $d\in\mathcal{O}$.

The key point is that the same $d$ works for every $\beta\ge b+1$. Indeed, if $g_\beta-dL(\beta)=\varpi y_\beta$ with
$y_\beta\in A_2(\beta)$ (by \ref{H:div}), then
$g_{\beta+1}-dL(\beta+1)=T_\beta(g_\beta-dL(\beta))=\varpi T_\beta y_\beta$, and $T_\beta y_\beta\in A_1(\beta+1)$ by \ref{H:T}.
Hence $L(\beta+2)-(k+\varpi^bd)L(\beta)=\varpi^{b+1}y_\beta$ with $y_\beta\in A_2(\beta)$, which is $\mathbf P(b+1,k+\varpi^bd)$.
The limit statement follows because $(K_j)$ is Cauchy and $L(b+2)-\kappa_\infty L(b)=(L(b+2)-K_{b-1}L(b))-(\kappa_\infty-K_{b-1})L(b)$.
\end{proof}

\begin{remark}
Only \ref{H:lead} uses $\dim S_{\ell-1}=1$. The original argument (docs, Step~3) instead showed that the leading digit
$c$ vanishes by comparison with Boylan--Webb stability at some level. The version above needs neither Boylan--Webb
stability nor the nonvanishing of $\Lb(b)$: the correction is absorbed into the constant. For $\ell=29,31$ everything
goes through except \ref{H:lead}, because $\dim S_{\ell-1}=2$. This matches the sharpness of $2m-1$ for those primes.
\end{remark}

\subsection{Assembly}\label{sec:assembly}
Let $M$ be the $\Zl$-module of cuspidal $q$-series having an integral Katz expansion. Let $T_b=D$ for $b$ even and
$T_b=U$ for $b$ odd, so that $T_bL(b)=L(b+1)$. Let $A_c(b)=A_c(\beta)$ for $b$ odd and $A_c(b)=A_c(\alpha)$ for $b$ even.
\begin{itemize}
\item \ref{H:L}: $L(0)=1$, and Lemma~\ref{lem:env}(2) applied inductively.
\item \ref{H:T}: Lemma~\ref{lem:env}(1).
\item \ref{H:lead}: For $x\in A_1(\gamma)$ we have $v_i(x)\ge1$ for $i\ge2$, so $x\equiv b_0+b_1\pmod\ell$ because $E\equiv1$. Here
  $b_0$ is a constant and must vanish modulo $\ell$ by cuspidality. Then $b_1\in\Sb_{\ell-1}=\Fl\Lb(b)$ by Lemma~\ref{lem:base}.
\item \ref{H:div}: uniqueness of Katz expansions modulo $\ell$, which is the mod-$\ell$ independence of the pieces
  $B_i$ (Swinnerton-Dyer~\cite{SwinnertonDyer1973}; formalized as \lean{ModP.katz\_independent}).
\item \ref{H:base}: Lemma~\ref{lem:base}.
\end{itemize}
Theorem~\ref{thm:induction} gives $\kappa_\ell:=\kappa_\infty$ with $L(b+2)\equiv\kappa_\ell L(b)\pmod{\ell^b}$ for all
$b\ge1$. Dividing by the common $\eta$-factor gives $P_\ell(b+2)\equiv\kappa_\ell P_\ell(b)\pmod{\ell^b}$. Comparing
coefficients gives~\eqref{eq:main} with $m=b$. \qed

\subsection{Translation to stabilization indices}\label{sec:translation}
\begin{proof}[Proof of Corollary~\ref{cor:bm}]
Fix $m$. For $b\ge m$, Theorem~\ref{thm:A} applied with $b$ in place of $m$ gives $L(b+2)\equiv\kappa_\ell L(b)\pmod{\ell^b}$,
hence modulo $\ell^m$. Since $\kappa_\ell$ is a unit, $L(b+2k)\equiv\kappa_\ell^kL(b)\pmod{\ell^m}$, so
$\Lambda(b,m)=(\Z/\ell^m)\Lb(b)=\Lambda(b+2,m)$. The module is free of rank one: by Lemma~\ref{lem:base} $L(b)\not\equiv0\pmod\ell$,
so $\Lambda(b,m)\cong\Z/\ell^m$. The map $X$ (or $U\circ D$ on the even chain) acts on it as multiplication by the unit
$\kappa_\ell$, and is therefore an isomorphism.
\end{proof}

% =====================================================================
\section{The Lean formalization}\label{sec:lean}
% =====================================================================

\subsection{The formal statement}
The theorem for $\ell=13$ reads as follows; the other three differ only in $\ell$ and $\kappa^{(0)}_\ell$. We render
Lean's Unicode in \LaTeX{} notation.
\begin{alltt}\small
theorem OrdQuot.atkin_obrien_13 :
    \(\exists\) \(\kappa\) : \(\mathbb{Z}\)_[13], (\(\exists\) d : \(\mathbb{Z}\)_[13], \(\kappa\) = 6 + 13 * d) \(\wedge\)
      \(\forall\) m, 1 \(\le\) m \(\to\) \(\forall\) n : \(\mathbb{N}\), 24 \(\mid\) 13 ^ m * n + 1 \(\to\)
        (13 : \(\mathbb{Z}\)_[13]) ^ m \(\mid\)
          (Fintype.card (Nat.Partition ((13 ^ (m + 2) * n + 1) / 24)) : \(\mathbb{Z}\)_[13])
          - \(\kappa\) * (Fintype.card (Nat.Partition ((13 ^ m * n + 1) / 24)) : \(\mathbb{Z}\)_[13])
\end{alltt}
Every constant in the statement comes from Lean core or Mathlib~\cite{Mathlib2020}: \lean{Nat.Partition} with Mathlib's
\lean{Fintype} instance, \lean{PadicInt}, divisibility in $\Zl$ and the casts. The one repository constant that appears
is an instance \lean{Fact (Nat.Prime 13)}. It is a proof of a true proposition and cannot change the meaning of the
statement. The division \lean{/} is truncating division on $\mathbb N$, but it is exact under the hypothesis
$24\mid13^mn+1$, and then also $24\mid13^{m+2}n+1$.

\subsection{From the paper proof to the formal proof}
The formal proof follows \S\S\ref{sec:setup}--\ref{sec:proofA}. It replaces rigid geometry by pointwise statements about
elliptic curves over $\Cl$:
\begin{itemize}
\item \emph{Katz series} are sequences of isobaric polynomials $b_i(E_4,E_6)$ in fixed complement pieces, with their
  $q$-expansions in $\Zl[[q]]$. Overconvergence enters only through Gauss-norm bounds on the $b_i$.
\item \emph{Test objects} are good-reduction Weierstrass curves over $\Cl$. A Katz series is evaluated at a curve through
  $(c_4,-c_6)$, and Proposition~\ref{prop:K} turns pointwise bounds into coefficient bounds.
\item The operators on Katz series are reached through classical Hecke theory. For $U$, the formalization uses
  $U(b/E^j)\equiv T_\ell(b\,E^{R-j})/E^R\pmod{\ell^N}$ with $R=(j+1)\ell^n$, and the geometric identity
  $\ell\,T_\ell F(W)=\sum_{C}F(W/C)$ over the $\ell+1$ subgroups $C$ of order $\ell$. In that sum the non-canonical terms
  satisfy the sharp bound and the canonical term is negligible.
\item The sharp bound (I2) is obtained from the $\Delta$-modular equation. The formalization uses Newton's identities for
  its coefficients, pointwise valuation bounds for the V\'elu quotients $\Delta(W/C)$, and a Hensel twelfth root in the
  sharp Katz class. Proposition~\ref{prop:K} is applied at radii $v(s)=(\ell M-1)/((\ell+1)M)$ with $M\to\infty$.
\item (DB) is derived from (UB) and (I2) through the identity
  $Df=Q^{1/2}\,U\big(Q^{\ell/2}G^{1/2}(\Delta^\delta/E^s)f\big)$, which follows from $\Phi=\Delta^\delta V(h)h^\ell$ with
  $h=\eta(\ell z)/\eta(z)^\ell$ and $h^2=Q/E$. The level-$\ell^2$ statement $V(Q)\in B_E(r)$ is therefore never needed.
\item $G=E/V(E)\in B_E(r)$ follows from the modular equation of $E_{\ell-1}$, whose coefficients are classical, by an
  $\ell$-adic Hensel contraction.
\item On the ordinary locus, good reduction of all quotients follows from $\prod_C\Delta(W/C)=\ell^{12}\Delta(W)^{\ell+1}$,
  the constant term of the $\Delta$-modular equation.
\end{itemize}

\subsection{Components built}
Table~\ref{tab:lean} lists the main components. Most of them are not in Mathlib. They are stated in reusable generality
(for example, Newton polygons over any ultrametric field, and Lubin's quotient over complete linearly topologized domains).
In total there are about $43{,}000$ lines in $38$ directories.

\begin{table}[ht]
\centering\small
\caption{Main components of the formal proof. Line counts are approximate.}\label{tab:lean}
\begin{tabular}{@{}p{0.27\textwidth}p{0.55\textwidth}r@{}}
\toprule
Directory & Content & Lines\\
\midrule
\lean{AtkinOBrien} & Envelope inequalities, abstract induction, Newton polygons over ultrametric fields, canonical subgroup (Newton half of Katz 3.10.7) & 1089\\
\lean{FormalGroupLayer} & Partial derivatives and the multivariable chain rule for power series; invariant differential; $[n]^*\omega=n\omega$; $[p]$ in characteristic $p$; height-$2$ coefficient ideals & 912\\
\lean{EllipticFormal} & Formal group of a Weierstrass curve over any commutative ring; invariant differential & 1539\\
\lean{LubinQuotient}, \lean{CanonicalIsogeny} & Lubin's quotient by a finite subgroup, the dual isogeny, $\|\Lie\check\pi\|=\|a\|$ & 2236\\
\lean{Velu}, \lean{VeluLubin}, \lean{VeluDescent} & V\'elu's formulas; formal V\'elu series; integral descent; link to Lubin & 5593\\
\lean{TateCurve}, \lean{TateFormal}, \lean{TateAdd}, \lean{TateQuot} & $\Delta(\mathrm{Tate})=q\prod(1-q^n)^{24}$; von Staudt--Clausen; Tate uniformization as a formal identity; the addition theorem; V\'elu quotients of the Tate curve & 6400\\
\lean{HasseInv} & Hasse invariant $\equiv E_{\ell-1}$ (Deligne's congruence) & 977\\
\lean{ModP}, \lean{Igusa}, \lean{KatzForms} & Swinnerton-Dyer structure mod $p$; Igusa's simple zeros; supersingular points; Step (K) & 2904\\
\lean{RamanujanODE} & $3\theta E_4=E_2E_4-E_6$, $2\theta E_6=E_2E_6-E_4^2$, $12\theta E_2=E_2^2-E_4$ over $\Z[[q]]$ & 637\\
\lean{HeightEll} & Height exactly $2$ for supersingular good reduction, $p\ge5$ & 1649\\
\lean{QuotientHasse}, \lean{AllQuotients}, \lean{QuotE}, \lean{OrdQuot} & Buzzard~3.3 for formal groups; quotient models and valuations for all $\ell+1$ subgroups; ordinary locus & 4416\\
\lean{HeckeQExp}, \lean{HeckeBridge}, \lean{SubgroupSum}, \lean{TorsionSpec} & $T_p$ on isobaric polynomials; geometric Hecke sum; torsion integrality; $\#E[n]=n^2$; specialization & 4355\\
\lean{UpContraction}, \lean{UBound}, \lean{UBSharp} & Vonk's Theorem~B; (UB) & 2246\\
\lean{VLayer}, \lean{SharpHQ}, \lean{ConjAlg}, \lean{DBound} & Modular equations; $G\in B_E(r)$; sharp (I2); (DB) & 4543\\
\lean{BoylanWebb}, \lean{AtkinOBrienFinal}, \lean{AtkinOBrienMain} & $L(b)$ as in~\cite{BoylanWebb2013}; certificates; Sturm bound; constant-absorption induction; assembly & 2996\\
\bottomrule
\end{tabular}
\end{table}

\subsection{Trust and reproduction}
The audit script \lean{scripts/audit/atkin\_obrien\_closure.lean} prints the constants in each statement and computes the
full transitive dependency closure of each proof. The closure has about $90{,}550$ constants, about $5{,}380$ of them
defined in the repository. Its only axioms are \lean{propext}, \lean{Quot.sound} and \lean{Classical.choice}, and
\lean{sorryAx} does not occur. A source scan of all $727$ Lean files found:
\begin{itemize}
\item no \lean{axiom} declarations;
\item no \lean{native\_decide} or \lean{ofReduceBool};
\item no \lean{unsafe}, \lean{implemented\_by}, \lean{extern}, \lean{csimp} or \lean{opaque};
\item no declaration-adding metaprograms and no kernel-skipping options.
\end{itemize}
Every intermediate hypothesis used during the development was eventually proved, since otherwise it would appear as an
argument of the final theorems. The hypotheses in question were the column bounds, (I2), (I3), $G$, the quotient
valuation bounds, the Tate kernels, subgroup-sum algebraicity, torsion integrality, specialization, Ramanujan's ODEs and
Deligne's congruence. Not yet done: an independent re-check of the compiled \lean{.olean} files with
\lean{lean4checker}.

\emph{Toolchain}: Lean \lean{v4.30.0-rc2}~\cite{Lean4}, Mathlib revision \lean{606efcb7}.
\emph{Reproduction}:
\begin{alltt}\small
git clone https://github.com/tomdif/RamanujanLean && cd RamanujanLean
lake exe cache get
lake build RamanujanTau.OrdQuot.Atkin
lake env lean scripts/audit/atkin_obrien_closure.lean
\end{alltt}
\TODO{Pin the repository commit (the theorem landed in \lean{5954d8d}); record build time and memory; run
\lean{lean4checker}; confirm that the default build target covers the audit imports.}

\begin{remark}[uniformity in $\ell$]
The envelope lemma is formalized for all $\ell\ge13$. The lemmas from which (UB), (I2) and (DB) are derived take a general
prime $\ell\ge5$ (respectively $\ell\ge11$ for (I2)). Only the base case (Lemma~\ref{lem:base}) and the final assembly are
specialized to the four primes. \TODO{Package (UB), (DB) and (I2) as standalone Lean theorems for all primes
$\ell\ge 13$. This would formally discharge the input G3 of \S\ref{sec:upper}.}
\end{remark}

\begin{remark}
The same repository contains formal proofs of Watson's theorem for powers of $5$ (\lean{ramanujan\_pow5}), the powers-of-$7$
theorem (\lean{ramanujan\_pow7}) and Atkin's theorem for powers of $11$ (\lean{ramanujan\_pow11\_unconditional}).
\TODO{Decide whether to mention these; if so, run the same closure audit on them.}
\end{remark}

% =====================================================================
\section{Proof of Theorem~\ref{thm:B}: $T_{\ell^2}$ is invertible modulo $\ell$}\label{sec:thmB}
% =====================================================================

Fix a prime $\ell\ge13$, with $j,t,\lambda$ as in \S1, and let
\[
  \rho=\frac{j\ell-1}{24},\qquad k'=\frac{\ell+j-2}{2}=\ell-1-12t,\qquad R=R_\ell .
\]
Then $\delta=\ell t+\rho$ and $\ell\rho+t=j\delta$. Put $\Sb=S_{\ell-1}(\Fl)$ and $\Mb=M_{\ell-1}(\Fl)=\Fl\cdot1\oplus\Sb$,
where $E_{\ell-1}\equiv1$, and
\[
  \cV:=\Delta^t\,S_{k'}(\Fl)\subseteq\Sb .
\]
Let $\Ub,\Db,\Xb$ be the reductions of $U,D,X$ on $q$-expansions modulo $\ell$. Following~\cite[(5.1)]{BoylanWebb2013}, the
Boylan--Webb invariant is
\begin{equation}\label{eq:dl}
  d_\ell=\min\{t'\ge0:\ f|D|X^{t'}\in\cV_{\mathrm{inv}}\ \text{for all } f\in M_{\ell-1}\},
\end{equation}
where $\cV_{\mathrm{inv}}$ is the largest subspace of $\Sb$ on which $\Xb$ is invertible.

\subsection{Mod-$\ell$ filtrations}
For a mod-$\ell$ modular form $f$ on $\Gamma_1(N)$, with $\ell\nmid N$ and $\ell\ge5$, let $\w(f)$ be its filtration. We use the
following standard facts~\cite{SerreBourbaki1973,SwinnertonDyer1973,Katz1977,Jochnowitz1982}:
\begin{enumerate}[label=(F\arabic*)]
\item $\w(\theta f)\le \w(f)+\ell+1$, with equality iff $\ell\nmid \w(f)$; in general
  $\w(\theta f)=\w(f)+\ell+1-s(\ell-1)$ with $s\ge0$.
\item For $k\ge2$, $U\equiv T_\ell$ on $M_k\bmod\ell$, and $\w(f|U)\le\ell+(\w(f)-1)/\ell$.
\item $f|U=0$ iff $\theta^{\ell-1}f=f$, since $f-\theta^{\ell-1}f=(f|U)|V$.
\item Filtrations of forms in one weight class are congruent modulo $\ell-1$. A nonzero form of filtration $0$ is constant,
  and $\theta$ kills constants.
\end{enumerate}
\TODO{Gap G2: cite the level-$N$ $\theta$-filtration theory at the precise statements (Katz 1977; Jochnowitz 1982, \S?;
Edixhoven 1992, \S3).}

\begin{lemma}[low point of a $\theta$-cycle]\label{lem:T}
Let $5\le k\le\ell-1$, and let $f\neq0$ be a mod-$\ell$ form with $\w(f)=k$ and $f|U=0$. Then $\w(\theta^{\ell+1-k}f)=\ell+3-k$.
\end{lemma}
\begin{proof}
Put $a_i=\w(\theta^if)$. By (F3), $a_{\ell-1}=a_0=k$, and by (F1) the total drop over the cycle is $\sum s_i=\ell+1$. Drops
occur only where $\ell\mid a_i$, and no $\theta^if$ vanishes. For $i\le\ell-k$ we have $a_i=k+i(\ell+1)$, so
$a_{\ell-k}=\ell(\ell+1-k)$ is divisible by $\ell$. Let $s$ be the first drop. Then $a_{\ell-k+1}=\ell(\ell+2-k)+1-s(\ell-1)\equiv1+s\pmod\ell$,
and after this the filtrations rise by $\ell+1$ per step until $i=2\ell-k-s$. Since the cycle closes at $i=\ell-1$, a second
drop is possible only if $s\ge\ell-k+2$. A single drop of size $\ell+1$ would make $a_{\ell-k+1}$ negative, so
$s\ge\ell-k+2$. Conversely, $s=\ell-k+3$ would give $a=4-k\le0$, and then $\theta^{\ell-k+1}f$ would be a constant in the image
of $\theta$, i.e.\ zero. Hence $s=\ell-k+2$ and $a_{\ell+1-k}=\ell+3-k$.
\end{proof}

\begin{lemma}\label{lem:Uinv}
For $\ell\ge7$, $\Ub=T_\ell$ is invertible on $\Sb$.
\end{lemma}
\begin{proof}
A nonzero $f\in\Sb$ has $\w(f)=\ell-1$, because the only smaller weight in its class is $0$. If $f|U=0$, Lemma~\ref{lem:T}
gives $\w(\theta^2f)=4$, so $\theta^2f=cE_4$. The constant terms force $c=0$, and then $f=\theta^{\ell-3}\theta^2f=0$.
\end{proof}

\begin{lemma}[image of $\Db$; the FKO bound]\label{lem:imageD}
For $\ell\ge13$, $\Db(\Mb)\subseteq\cV$. In particular $\dim\Db(\Mb)\le\dim\cV=R$.
\end{lemma}
\begin{proof}
Modulo $\ell$ we have $\Phi\equiv\Delta^\delta=V(\Delta^t)\Delta^\rho$, so $\Db f=\Delta^t\,U(\Delta^\rho f)$. The form
$\Delta^\rho f$ has weight $12\rho+\ell-1$, and by (F2)
$\w(U(\Delta^\rho f))\le\ell+(12\rho+\ell-2)/\ell<\ell+(j+2)/2$. The class of this filtration modulo $\ell-1$ is
$12\rho\equiv k'$. The next weight in the class, $k'+\ell-1=(3\ell+j-4)/2$, exceeds $\ell+(j+2)/2$ when $\ell>6$, so the
filtration is at most $k'$. Finally, $\ord_\infty(\Delta^\rho f)\ge\rho\ge1$, so $U(\Delta^\rho f)$ has no constant term and
lies in $S_{k'}$. This also covers $f=1$.
\end{proof}

\begin{lemma}\label{lem:dims}
For $\ell\ge13$, $\dim S_{k'}=R=\lceil d/2\rceil$, where $d=\dim S_{\ell-1}$.
\end{lemma}
\begin{proof}
Check each $j$. If $j\le11$, then $d=2t$ and $\dim S_{12t+j-1}=t=\lfloor(\ell+12)/24\rfloor$. If $j\ge13$, then $d=2t+1$
and $\dim S_{12t+j-1}=t+1=\lfloor(\ell+12)/24\rfloor$.
\end{proof}

\subsection{The half-integral-weight dictionary}\label{sec:dict}
Recall $\cS=\cS_\ell=S_{\lambda+1/2}(\SL,\nu_\eta^{-\ell})$; note $\nu_\eta^{-\ell}=\nu_\eta^{-j}$.
\begin{enumerate}[label=(D\arabic*)]
\item Division by $\eta^j$ is an isomorphism $S_{k'}\to\cS$: the weight is $k'-j/2=\lambda+\frac12$, and $G\in S_{k'}$ has
  $\ord_\infty G\ge1>j/24$. Since $\eta^\ell=\Delta^t\eta^j$, we get $\cV=\eta^\ell\cS\equiv\eta(\ell z)\cS\pmod\ell$.
\item \emph{$\Xb|_\cV$ is $U_{\ell^2}$.} Let $F=\Delta^tG\in\cV$. Using $U(V(A)B)=A\,U(B)$ and $V(\Delta^\rho)\equiv\Delta^{\ell\rho}$, we get
  $\Db\Ub F=\Delta^t\,U^2(\Delta^{\ell\rho+t}G)=\Delta^t\,U^2(\Delta^{j\delta}G)$. Since
  $\Delta^{j\delta}\equiv V_{\ell^2}(\eta^j)/\eta^j$, this equals $\Delta^t\eta^j\,U_{\ell^2}(G/\eta^j)$. So in the coordinate
  $h=G/\eta^j\in\cS$, $\Xb|_\cV=U_{\ell^2}$. In particular $\Xb(\Sb)\subseteq\cV$ and $\Xb(\cV)\subseteq\cV$.
\item \emph{$U_{\ell^2}\equiv T_{\ell^2}$.} On $h(24z)\in S_{\lambda+1/2}(\Gamma_0(576),\chi)$,
  $T_{\ell^2}=U_{\ell^2}+\chi^*(\ell)\ell^{\lambda-1}\leg{\cdot}{\ell}+\ell^{2\lambda-1}V_{\ell^2}$, which is congruent to $U_{\ell^2}$
  modulo $\ell$ for $\lambda\ge2$. Moreover $T_{\ell^2}$ is integral, and $\cS(\Z)\otimes\Fl=\cS(\Fl)$ via the Miller basis of $S_{k'}$.
\item \emph{Odd/even halves.} Write $\Sb=\eta\cdot\cT$ with $\cT=S_{\ell-3/2}(1,\nu_\eta^{-1})$. Then $\Db(\eta g)=\eta(\ell z)U_\ell(g)$
  holds exactly, because $\Phi\eta=\eta(\ell^2z)$, and $\Ub(\eta(\ell z)h)=\eta\,U_\ell(h)$. So $\Db$ and $\Ub|_\cV$ are the two
  halves $\cT\to\cS$ and $\cS\to\cT$ of $U_{\ell^2}=U_\ell\circ U_\ell$.
\item \emph{The partition vector.} $\Lb(1)=\Db(1)=\eta(\ell z)U_\ell(1/\eta)$. In the coordinates
  $W(G)=$ (the coefficients of $q^1,\dots,q^R$ in $G\prod_n(1-q^n)^{-j}$), it is the vector $(p(\ell n-\rho)\bmod\ell)_{1\le n\le R}$.
\end{enumerate}

\subsection{Yang's isomorphism}
\begin{theorem}[Yang~{\cite[Thm.~1]{Yang2014}}]\label{thm:yang}
Let $0<r<24$ with $(r,6)=1$, and let $s\ge0$ be even. Put $S_{r,s}=\{\eta(24\tau)^rf(24\tau):f\in M_s(1)\}\subset S_{r/2+s}(576,\leg{12}{\cdot})$.
The Shimura correspondence gives an isomorphism of Hecke modules
\[
  S_{r,s}\ \cong\ S^{\mathrm{new}}_{r+2s-1}\big(6,\ -\leg{8}{r},\ -\leg{12}{r}\big)\otimes\leg{12}{\cdot},
\]
where the last two entries are the eigenvalues of the Atkin--Lehner involutions $W_2$ and $W_3$. Under this isomorphism
$T_{p^2}$ corresponds to $\leg{12}{p}T_p$ for $p\nmid6$.
\end{theorem}

\begin{remark}[checking the citation]\label{rem:yangcheck}
We apply the theorem with $r=24-j$ and $s=12t+j-13=(\ell-r-2)/2$. Then:
\begin{itemize}
\item $\nu_\eta^{-\ell}=\nu_\eta^r$, and $S_{r,s}(1)=\eta^rM_s(1)$~\cite[Prop.~6]{Yang2014} is all of $\cS$; this agrees
  with (D1) because $G/\eta^j=\eta^r\cdot G/\Delta$.
\item The weights match: $(r-1)/2+s=\lambda$ and $r+2s-1=\ell-3$.
\item The hypotheses hold for every $\ell\ge13$: $s$ is even, and $s\ge0$ exactly when $\ell\ge13$.
\item The characters $\leg{8}{\cdot}$ and $\leg{12}{\cdot}$ are even, so the sign class is $(-\leg{2}{\ell},-\leg{3}{\ell})$.
\item The Atkin--Lehner convention is $a_p=-\varepsilon_p\,p^{k/2-1}$, which is the same as PARI's.
\item Yang's $T_{\ell^2}$~\cite[Prop.~11]{Yang2014} is the operator in (D3).
\item The map is onto the whole newform eigenspace of the class, and multiplicity one holds~\cite[Cor.~1]{Yang2014}.
\end{itemize}
These points were checked against arXiv:1110.1810v1, and independently by an external reviewer. The published version
was not opened. Numerically, for $\ell\in\{29,31,37,41,43,47,53,59,61,67,71,73,83\}$, which covers all eight classes of
$\ell\bmod24$, the class $(-\leg2\ell,-\leg3\ell)$ is the one whose $T_\ell$-characteristic polynomial modulo $\ell$, after
$x\mapsto\leg{12}{\ell}x$, equals $\charpoly(\Xb|_\cV)$. A separate PARI scan confirmed $a_2=-w_22^{(\ell-5)/2}$ and
$a_3=-w_33^{(\ell-5)/2}$ for every newform in $S^{\mathrm{new}}_{\ell-3}(\Gamma_0(6))$ for the primes $13\le\ell\le103$
other than $79$ and $89$. \TODO{Check against the published version; complete the scan for $\ell=79,89$.}
\end{remark}

Consequently $\det(T_{\ell^2}|\cS)=\leg{12}{\ell}^R\prod_Fa_\ell(F)$, where $F$ runs over the newforms of weight $\ell-3$ on
$\Gamma_0(6)$ in the class $(-\leg2\ell,-\leg3\ell)$, counted with all embeddings.

\subsection{The proof}
\begin{proof}[Proof of Theorem~\ref{thm:B}]
Suppose some newform $F$ in the class has $a_\ell(F)\equiv0\pmod{\mathfrak l}$ for a prime $\mathfrak l\mid\ell$. Its reduction
$\bar F$ is a nonzero mod-$\mathfrak l$ eigenform on $\Gamma_0(6)$ with $\bar F|U=0$. Its filtration is $\ell-3$, since the only
smaller weight in its class is $-2$. Lemma~\ref{lem:T} with $k=\ell-3$ gives $\w(\theta^4\bar F)=6$. So $\tilde G:=\theta^4\bar F$ is a
nonzero mod-$\mathfrak l$ cuspidal eigenform in $S_6(\Gamma_0(6))\otimes\bar\F_\ell$. Since $a_{pn}(\theta^4F)=p^4n^4a_{pn}(F)$, its
eigenvalues are $p^4a_p(F)$ for $T_p$ with $p\nmid6\ell$, $16a_2(F)$ for $U_2$, and $81a_3(F)$ for $U_3$. By the Deligne--Serre lifting
lemma~\cite{DeligneSerre1974}, this system of eigenvalues lifts to an eigenform in $S_6(\Gamma_0(6))$.

The space $S_6(\Gamma_0(6))$ is three-dimensional, because $S_6(\Gamma_0(2))=S_6(\SL)=0$ and $S_6(\Gamma_0(3))$ is one-dimensional. It is spanned by:
\begin{itemize}
\item the newform $6.6.a.a=q+4q^2-9q^3+16q^4-66q^5+\cdots$, with $(w_2,w_3)=(-1,+1)$;
\item the two $U_2$-eigen-oldforms from $3.6.a.a$, whose $U_2$-eigenvalues are the roots of $x^2+6x+32$ and whose $U_3$-eigenvalue is $9$.
\end{itemize}
These data were computed with PARI/GP~\cite{PARI}; see also~\cite{LMFDB}.

Next we compute the $U_2$ and $U_3$ eigenvalues of $\tilde G$. Since $w_2(F)=-\leg2\ell$, we have
$a_2(F)=\leg2\ell2^{(\ell-5)/2}$. Because $2^{(\ell-1)/2}\equiv\leg2\ell$, this gives $a_2(F)\equiv2^{-2}$; likewise
$a_3(F)\equiv3^{-2}$. Hence $\tilde G$ has $U_2$-eigenvalue $4$ and $U_3$-eigenvalue $9$. This rules out both
possibilities for the lift:
\begin{itemize}
\item The newform $6.6.a.a$ has $U_3$-eigenvalue $-9$. Then $9\equiv-9$ would force $\ell\mid18$, which is impossible.
\item For the oldforms, $U_2$-eigenvalue $4$ would require $16+24+32=72\equiv0\pmod\ell$, which is impossible for $\ell\ge5$.
\end{itemize}
Therefore every $a_\ell(F)$ is a unit, and $\det(T_{\ell^2}|\cS)$ is an $\ell$-adic unit.
\end{proof}

\begin{remark}[why the sign matters; $\ell=859$]\label{rem:859}
The newform $6.6.a.a$ is not ordinary at $\ell=859$: $a_{859}=-33566284=-859\cdot39076$. This was found by a PARI scan of
$\ell\le3\cdot10^5$, which was interrupted after $859$; we have not established whether $859$ is the smallest such prime.
For this $\ell$, $\theta^{\ell-5}(6.6.a.a)$ is a non-ordinary eigenform of weight $\ell-3$ in the class
$(-\leg2\ell,+\leg3\ell)$, which is the \emph{wrong} class. If Yang's class had the sign $+\leg3\ell$, Theorem~\ref{thm:B} would fail at
$\ell=859$, contradicting Boylan--Webb's computation $d_{859}=0$. A direct computation at $\ell=859$ confirms
Theorem~\ref{thm:B} (\S\ref{sec:numerics}): $\det(\Xb|_\cV)\equiv318\not\equiv0\pmod{859}$.
\end{remark}

\subsection{Proof of Corollary~\ref{cor:dl}}
(i) By Theorem~\ref{thm:B} and (D2), $\cV=\Xb(\cV)\subseteq\Db(\Sb)\subseteq\Db(\Mb)\subseteq\cV$, using
Lemma~\ref{lem:imageD}, so all of these are equal. Lemma~\ref{lem:dims} gives the rank.
(ii) $\ker\Xb=\Ub^{-1}(\ker\Db)$ has dimension $d-R$ by Lemma~\ref{lem:Uinv}. Since $\Xb(\Sb)\subseteq\cV$ and
$\Xb|_\cV$ is invertible, $\Sb=\cV\oplus\ker\Xb$ and $\rank\Xb^k=R$ for all $k\ge1$.
(iii) By (ii), $\cV_{\mathrm{inv}}=\operatorname{im}\Xb^d=\cV\supseteq\Db(\Mb)$, so $d_\ell=0$ in~\eqref{eq:dl}.
(iv) This is~\cite[Thm.~1.2]{BoylanWebb2013}.
(v) The number of zero slopes of $X$ on the odd overconvergent space equals the dimension of the generalized nonzero
eigenspace of $\Xb$ on the reduction of an $X$-stable lattice. On the envelope lattice $A_1$, $\Xb$ factors through the
reduction $A_1\to\Sb$ to block $1$, whose nonzero generalized eigenspace is $\cV$, of dimension $R$. \qed

\begin{remark}
Note that $d_\ell=0$ is a condition on all of $M_{\ell-1}$, including the Eisenstein direction. An exact computational test is
$d_\ell=0\iff\rank[\Xb^d\mid\Db|_{\Mb}]=\rank\Xb^d$, which needs no semisimplicity assumption.
\end{remark}

\subsection{An independent proof of the upper bound}\label{sec:upper}
Assume the envelope lemma (Lemma~\ref{lem:env}) for $\ell$. Its inputs (I1)--(I3) are standard, but this is gap G3, and
it is formalized only in the sense of the uniformity remark in \S\ref{sec:lean}. Then Theorem~\ref{thm:B} gives
$b_\ell(m)\le2m-1$ without Boylan--Webb's $d_\ell$ machinery.

Let $Y\oplus Z$ be the Riesz decomposition of $X$ on the odd space, with $Y$ the unit-root part. Write $L(1)=y+z$ with
$y=e_YL(1)$, where $e_Y=\lim X^{n!}$. The reduction $\rho:A_1\to\Sb$ has kernel $\ell A_2$ and is $X$-equivariant. If
$x\in Z\cap A_1$, then $\rho(x)$ lies in the generalized kernel of $\Xb$, which is $\ker\Xb$ by Corollary~\ref{cor:dl}(ii). By
induction, $w_k:=X^kz/\ell^{k+1}\in A_2\cap Z$ for all $k$:
\begin{itemize}
\item \emph{Base.} $\rho(z)\in\ker\Xb\cap\cV=0$.
\item \emph{Step.} $\rho(Xw_k)=\Db\rho_\alpha(Uw_k)\in\cV\cap\ker\Xb=0$.
\end{itemize}
So $v(X^kz)\ge k+1$, which gives $b_\ell(m)\le2m-1$ (Proposition~\ref{prop:reduction}). For $d=1$ we have $\ker\Xb=0$, and the
same argument run one $T$-step at a time is the proof of Theorem~\ref{thm:A}.

% =====================================================================
\section{Rank attainment and the prime $607$}\label{sec:rank}
% =====================================================================

Assume Theorem~\ref{thm:B}. Then $Y$ has rank $R$, $\bar z=0$, and $\bar y=\Lb(1)=\eta(\ell z)U_\ell(1/\eta)\in\cV$. The stable
odd module is $\Omega(m)=\Zl[X]y\bmod\ell^m$. Since $\Xb|_\cV$ is invertible, the mod-$\ell$ stable span of the $\Lb(b)$ is
the Krylov span $\Fl[\Xb]\Lb(1)$. Hence:

\begin{proposition}\label{prop:krylov}
$r_\ell(1)=\dim_{\Fl}\Fl[T_{\ell^2}]\cdot U_\ell(1/\eta)$, the Krylov rank of $\Lb(1)$ in $\cS\otimes\Fl\cong\cV$.
\end{proposition}

\begin{proposition}[Nakayama direction]\label{prop:nakayama}
If $\Lb(1)$ is a cyclic vector of $\Xb|_\cV\cong T_{\ell^2}|_{\cS\otimes\Fl}$, then $\Omega(m)$ is free of rank $R$ for every $m\ge1$.
\end{proposition}
\begin{proof}
$\Zl[X]y$ is a submodule of the rank-$R$ lattice of $Y$ whose image modulo $\ell$ is all of $\bar Y=\cV$. By Nakayama's lemma it is the
whole lattice.
\end{proof}

\begin{remark}[gap G6]\label{rem:G6}
The converse fails at the level of linear algebra: non-cyclicity modulo $\ell$ does not by itself force a rank deficit at
higher precision. For example, $A=\mathrm{diag}(1,1+\ell)$ and $y=(1,1)$ generate a rank-$2$ lattice whose image modulo $\ell$ is
one-dimensional. So the statement ``$r_\ell(m)$ equals the Krylov rank of $\Lb(1)$ for all $m$'' needs a lattice argument
that we have not written out. In the same vein, cyclicity of $\Lb(1)$ requires that the minimal polynomial of $T_{\ell^2}$
modulo $\ell$ equal its characteristic polynomial, and that $\Lb(1)$ have a nonzero component in each primary component.
Repeated eigenvalues are neither necessary nor sufficient for non-cyclicity. A geometric multiplicity $\ge2$ is
sufficient.
\end{remark}
\TODO{Re-read the end of docs/BL\_SHARPNESS\_PROOF.md before submission: a gap-closing pass on G6 is in progress. On the
date of this draft no ``Gap closure'' section existed.}

\subsection{The prime $\ell=607$}\label{sec:607}
\st{Computed}. The computation used \texttt{halfint.py} over $\F_{607}$, with $q$-expansions to $\ell^2R\approx9.2\cdot10^6$ terms
(about 30 minutes). It gave $d=50$ and $R=25$. Moreover:
\begin{itemize}
\item $\rank\Xb|_\cV=25$ and $\det\neq0$. So Theorem~\ref{thm:B} holds at $607$: $d_{607}=0$, $\rank\Db=25$, and $\Xb$ is
  semisimple at $0$.
\item $\charpoly(\Xb|_\cV)=f_{22}(x)\,(x+201)\,(x+367)^2$, where $f_{22}$ is irreducible of degree $22$.
\item $\dim\ker(\Xb|_\cV+367)=2$. So $T_{607^2}$ has a two-dimensional eigenspace modulo $607$, and no vector of
  $\cS\otimes\F_{607}$ is cyclic.
\item The Krylov rank of $\Lb(1)$ is $24=R-1$. This equals Boylan--Webb's $r_{607}=R_{607}-1$~\cite[\S6.1]{BoylanWebb2013}.
\end{itemize}
Under Yang's isomorphism, the repeated eigenvalue means that two Galois orbits of newforms of weight $604$ on $\Gamma_0(6)$ in the
class have congruent $a_{607}$ modulo a prime above $607$. The deficit at $607$ is therefore forced by the Hecke module
itself, not by the particular vector, and it is consistent with $d_{607}=0$. Boylan--Webb assert only that $d_\ell>0$
forces a deficit, not the converse~\cite[\S6.2]{BoylanWebb2013}. The composition method of \S\ref{sec:methods} reproduces the
factor $(x+367)^2$ exactly.
\TODO{Determine whether Boylan--Webb prove $r_{607}(m)=R-1$ for all $m$ or compute it for $m=1$ only.}
\TODO{``Two Galois orbits'' versus ``two embeddings of one orbit'': determine which occurs.}

\subsection{A heuristic for the exceptional set}
A random monic polynomial over $\Fl$ fails to be squarefree with probability $1/\ell$. A fixed vector fails to be cyclic for
a regular matrix with probability about $(\#\text{primary components})/\ell^{\deg}$, which is $\lesssim1/\ell$. The model
therefore predicts that rank attainment fails at $\ell$ with probability $\asymp1/\ell$. Summing over primes, it predicts
about $\sum_{29\le\ell\le1300}1/\ell\approx0.8$ failures below $1300$ (Boylan--Webb found exactly one, $\ell=607$), and
infinitely many failures in total, forming a set of relative density zero among the primes. \st{Heuristic.}

% =====================================================================
\section{Sharpness of $2m-1$: reduction and numerics}\label{sec:numerics}
% =====================================================================

\subsection{Reduction to valuations}
\begin{proposition}[assuming Theorem~\ref{thm:B}]\label{prop:reduction}
With $z=L(1)-e_YL(1)$ as in \S\ref{sec:upper},
\[
  b_\ell(m)=\min(2k_o+1,\,2k_e+2),\quad k_o=\min\{k:v(X^kz)\ge m\},\quad k_e=\min\{k:v(UX^kz)\ge m\}.
\]
Hence $b_\ell(m)\le2m-1$ for all $m$ iff $v(X^kz)\ge k+1$ for all $k$, and $b_\ell(m)=2m-1$ for all $m$ iff in addition
$v(UX^kz)=k+1$ for all $k$.
\end{proposition}
\begin{proof}[Sketch]
$q$-expansion integrality is equivalent to integrality of Katz coordinates. Since $X|_Y$ is invertible and $X^{-1}|_Y$ is a
polynomial in $X$, $\Span\{X^iy:i\ge k\}=\Zl[X]y$; this gives ``$\Leftarrow$''. For ``$\Rightarrow$'', project with the
integral idempotent $e_Z$.
\end{proof}

\subsection{The second-order operator $\Theta$}\label{sec:theta}
Let $u_k=Uw_k\in A_1(\alpha)$ and $\bar u_k=\rho(u_k)\in\Sb$. By \S\ref{sec:upper}, $\Db\bar u_k=0$, so
$\bar u_k\in\bar K:=\ker\Db|_{\Sb}$, which has dimension $\lfloor d/2\rfloor$. Write $\Sb=\Ub(\cV)\oplus\bar K$, and let $\pi_K$ be
the projection along $\Ub(\cV)$. Define
\[
  \Theta:\bar K\to\bar K,\qquad\Theta(\bar u)=\pi_K\rho\big(U(D\tilde u/\ell)\big),
\]
for any lift $\tilde u\in A_1(\alpha)$ of $\bar u$. This is well defined: $D\tilde u/\ell\in A_2$, and changing the lift changes the
value by an element of $\Ub\Db(\cdot)\subseteq\Ub(\cV)$. Then $\bar u_{k+1}=\Theta\bar u_k$, so
\[
  b_\ell(m)=2m-1\ \text{for all }m\iff\Theta^k\bar u_0\ne0\ \text{for all }k,
\]
and a sufficient condition is that $\Theta$ is invertible and $\bar u_0\ne0$. In half-integral terms,
$\Theta(\bar u)=\pi_K\big[(U_{\ell^2}\tilde g)/\ell-U_{\ell^2}(e_1\tilde g)\big]$, where $\tilde u=\eta\tilde g$ and $E_{\ell-1}=1+\ell e_1$.
The first term is the slope-$1$ part of $T_{\ell^2}$ on the classical weight-$(\ell-\frac32)$ space; the second is a
weight-derivative correction. \st{Proved for the reduction} (given Theorem~\ref{thm:B} and the envelope lemma);
\st{open} in general.

\begin{computation}[conditional]\label{comp:theta}
For every prime $29\le\ell\le113$, $\Theta$ has full rank $\lfloor d/2\rfloor$ with squarefree characteristic polynomial, and
$\bar u_0\ne0$ (Table~\ref{tab:theta}). Hence $b_\ell(m)=2m-1$ for \emph{all} $m\ge1$ for these $\ell$, subject to three
conditions: (a) Theorem~\ref{thm:B}; (b) the envelope lemma for these $\ell$ (gap G3); and (c) the validity of the truncated Katz
model ($\mathrm{PR}=3$, $N=6$) used to compute $\Theta$. Condition (c) holds because $\Theta$ modulo $\ell$ sees only Katz
blocks $i$ with $\alpha(i),\beta(i)\le2$, but this routine justification has \emph{not yet been written out}. The
independent external review reproduced $\Theta$ for all $21$ primes and recorded the same qualification.
\end{computation}

For $29\le\ell\le47$ we have $\dim\bar K=1$, and the eigenvalue of $\Theta$ equals $\mu=\lambda_{\text{slope }1}/\ell\bmod\ell$, read off
from the Newton polygon of $\charpoly(X)$ in the Katz model. This supports the expectation that $\Theta$ is the reduction of $X/\ell$
on the slope-$1$ part.

\begin{table}[ht]
\centering\small
\caption{$\Theta$ on $\bar K$ (\texttt{theta2.py}, Katz model $\mathrm{PR}=3$, $N=6$). For $\dim\bar K=1$ we list the eigenvalue $\mu$,
$\bar u_0$ and $\bar u_1=\Theta\bar u_0$; for $\dim\bar K\ge2$, the factorization of $\charpoly\Theta$ over $\Fl$.}\label{tab:theta}
\begin{tabular}{@{}rcl@{}}
\toprule
$\ell$ & $\dim\bar K$ & $\Theta$\\
\midrule
29 & 1 & $\mu=23$, $\bar u_0=8$, $\bar u_1=10$\\
31 & 1 & $\mu=2$, $\bar u_0=5$, $\bar u_1=10$\\
37 & 1 & $\mu=20$, $\bar u_0=9$, $\bar u_1=32$\\
41 & 1 & $\mu=31$, $\bar u_0=26$, $\bar u_1=27$\\
43 & 1 & $\mu=31$, $\bar u_0=40$, $\bar u_1=36$\\
47 & 1 & $\mu=17$, $\bar u_0=19$, $\bar u_1=41$\\
53 & 2 & $(x+18)(x+28)$\\
59 & 2 & $(x+46)(x+58)$\\
61 & 2 & $x^2+44x+4$\\
67 & 2 & $x^2+35x+17$\\
71 & 2 & $x^2+34x+30$\\
73 & 3 & $(x+14)(x^2+31x+43)$\\
79 & 3 & $x^3+31x^2+22x+32$\\
83 & 3 & $(x+53)(x^2+64x+57)$\\
89 & 3 & $x^3+41x^2+27x+81$\\
97 & 4 & $(x+82)(x^3+36x^2+70x+27)$\\
101 & 4 & $(x+5)(x+3)(x^2+92x+93)$\\
103 & 4 & $(x+64)(x^3+41x^2+14x+56)$\\
107 & 4 & $x^4+62x^3+81x+84$\\
109 & 4 & $(x+81)(x+74)(x^2+23x+81)$\\
113 & 4 & $x^4+51x^3+17x^2+94x+109$\\
\bottomrule
\end{tabular}
\end{table}

\subsection{Computational methods}\label{sec:methods}
\begin{enumerate}[label=(M\arabic*)]
\item \emph{Katz-coordinate iteration} (\texttt{katzmat.py}, \texttt{bstab.py}, \texttt{bigstab.py}). We iterate $L(b)$ in Katz coordinates
  using the level-$1$ $U$- and $D$-matrices, and compute the stable modules exactly as Krylov spans, with Smith invariants over
  $\Z/\ell^m$. The model was validated in three ways:
  \begin{itemize}
  \item truncation stability: $N=\mathrm{PR}+3$ and $N=\mathrm{PR}+6$ agree in every case;
  \item zero mismatches against $p(n)$ data for $\ell=17,23,29,31$;
  \item it reproduces the classical cases $\ell=5,11$ ($b=m$), $\ell=7$ (Watson's $7^{\lfloor b/2\rfloor+1}$) and $\ell=13$
    (Atkin--O'Brien), and the values $\kappa^{(0)}_\ell=6,6,10,5$.
  \end{itemize}
\item \emph{Partition numbers.} We computed $p(n)$ modulo $17^5$, $19^4$, $23^4$ (and modulo powers of $13$, $29$, $31$) for $n\le10^8$
  by FLINT series inversion~\cite{FLINT}.
\item \emph{Half-integral weight} (\texttt{halfint.py}). This computes $\Xb|_\cV$ as $U_{\ell^2}$ in the coordinates $W$ of (D5), using
  $q$-expansions to $\ell^2R$ terms. For every $29\le\ell\le113$, its characteristic polynomial agrees factor by factor with the nonzero
  part of $\charpoly(\Xb|_{\Sb})$ computed independently in Katz coordinates.
\item \emph{Composition} (\texttt{l859\_compose.py}). This composes $\Db\circ\Ub$ with re-expansion in the basis $\Delta^iE_4^aE_6^b$. It
  needs only $O(\ell d)$ $q$-coefficients: at $\ell=859$, about $7\cdot10^4$ terms, $7$ seconds and $280$ MB. It was validated at
  $\ell=29,47,71,113$ and $607$. Before the run, an external reviewer checked it against the original script's full matrices at
  $15$ smaller primes and recomputed $\ell=859$ by the $\ell^2R$ method (about $2.6\cdot10^7$ coefficients, $4.5$ minutes, $3.5$ GB).
\item \emph{Level~$6$} (\texttt{yang.gp}). PARI/GP computation of $T_\ell$ modulo $\ell$ on $S^{\mathrm{new}}_{\ell-3}(\Gamma_0(6))$, split by
  Atkin--Lehner signs.
\end{enumerate}

\subsection{Data}
\begin{table}[ht]
\centering\small
\caption{Constants $\kappa_\ell\bmod\ell^m$ determined from $p(n)$, $n\le10^8$ (method M2), with the number of independent tests in
parentheses where recorded. In each case the proportionality fails modulo $\ell^{m+1}$ at $b=m$. The entries are compatible:
each reduces to the previous one, consistent with a single $\kappa_\ell\in\Zl$.}\label{tab:kappa}
\begin{tabular}{@{}rllll@{}}
\toprule
$\ell$ & $m=1$ & $m=2$ & $m=3$ & $m=4$\\
\midrule
13 & 6 & 45 (Atkin--O'Brien's $K_2$) & 214 & --\\
17 & 6 (20354) & 210 (1197) & 3678 (71) & 38069 (4)\\
19 & 10 & 48 & 3658 (41) & --\\
23 & 5 & 465 & 4697 (16) & --\\
\bottomrule
\end{tabular}
\end{table}

\begin{table}[ht]
\centering\small
\caption{Katz-coordinate sweep (method M1). The columns give: $d=\dim S_{\ell-1}$; the FKO bound $R_\ell$; the largest $m$
tested; the observed $b_\ell(m)$ for all tested $m$; and the stable rank modulo $\ell^m$ (the same for all tested $m$). The
stable modules are free in every case, truncation stability holds, and the unit-root count equals $R_\ell$. Rows with
$m_{\max}=16$ come from \texttt{bigm.log}; the others from \texttt{scan\_5\_60.json} and \texttt{scan\_60\_230.json}.}\label{tab:sweep}
\begin{tabular}{@{}lcccc@{}}
\toprule
$\ell$ & $d$ & $R_\ell$ & $m_{\max}$ & $b_\ell(m)$, rank\\
\midrule
13, 17, 19, 23 & 1 & 1 & 16 & $m$, rank 1\\
29, 31 & 2 & 1 & 16 & $2m-1$, rank 1\\
37, 41 & 3 & 2 & 16 & $2m-1$, rank 2\\
43, 47 & 3 & 2 & 9 & $2m-1$, rank 2\\
53, 59 & 4 & 2 & 9 & $2m-1$, rank 2\\
61, 67, 71 & 5 & 3 & 9 & $2m-1$, rank 3\\
73, 79, 83 & 6 & 3 & 9 & $2m-1$, rank 3\\
89 & 7 & 4 & 9 & $2m-1$, rank 4\\
97, 101, 103, 107 & 8 & 4 & 9 & $2m-1$, rank 4\\
109, 113 & 9 & 5 & 9 & $2m-1$, rank 5\\
127, 131 & 10 & 5 & 8 & $2m-1$, rank 5\\
137, 139 & 11 & 6 & 8 & $2m-1$, rank 6\\
149, 151 & 12 & 6 & 8 & $2m-1$, rank 6\\
157, 163, 167 & 13 & 7 & 8 & $2m-1$, rank 7\\
173, 179 & 14 & 7 & 8 & $2m-1$, rank 7\\
181, 191 & 15 & 8 & 8 & $2m-1$, rank 8\\
193, 197, 199 & 16 & 8 & 8 & $2m-1$, rank 8\\
211 & 17 & 9 & 8 & $2m-1$, rank 9\\
223, 227 & 18 & 9 & 7 & $2m-1$, rank 9\\
229 & 19 & 10 & 7 & $2m-1$, rank 10\\
\bottomrule
\end{tabular}
\end{table}
For the primes $5,7,11$ the same sweep gives $b_\ell(m)=m$ for $\ell=5,11$, and $1,2,4,6,\dots,2m-2$ for $\ell=7$, consistent with
Watson and Atkin. For $29\le\ell\le179$ the Newton polygon of $X$ in the Katz model shows $R_\ell$ slopes $0$ followed by
exactly $\lfloor d/2\rfloor$ slopes $1$; for $d=1$ the slopes are $0,2,2,4,\dots$ (no slope $1$). For $181\le\ell\le229$ the
truncated models ($\mathrm{PR}=8,7$) show fewer than $\lfloor d/2\rfloor$ slopes $1$. We attribute this to truncation, but
we have not checked it. \TODO{Recompute the slopes for $181\le\ell\le229$ at higher precision.}

\begin{table}[ht]
\centering\small
\caption{$\charpoly(\Xb|_{\Sb_{\ell-1}})$ over $\Fl$ for $13\le\ell\le113$ (\texttt{structure\_13\_60.log},
\texttt{structure\_61\_230.log}). In every case $\rank\Ub=d$, $\rank\Db=\rank\Xb=\rank\Xb^k=R_\ell$, and the factor $x^{d-R}$ is
the semisimple kernel of Corollary~\ref{cor:dl}(ii).}\label{tab:charpoly}
\begin{tabular}{@{}rl@{}}
\toprule
$\ell$ & $\charpoly(\Xb|_{\Sb})$\\
\midrule
13 & $x+7$\\
17 & $x+11$\\
19 & $x+9$\\
23 & $x+18$\\
29 & $(x+26)\,x$\\
31 & $(x+1)\,x$\\
37 & $(x+22)(x+3)\,x$\\
41 & $(x+18)(x+10)\,x$\\
43 & $(x+12)(x+37)\,x$\\
47 & $(x^2+18x+19)\,x$\\
53 & $(x+46)(x+10)\,x^2$\\
59 & $(x^2+53x+57)\,x^2$\\
61 & $(x+29)(x^2+2x+18)\,x^2$\\
67 & $(x+26)(x^2+42x+65)\,x^2$\\
71 & $(x^3+18x^2+2x+2)\,x^2$\\
73 & $(x^3+65x^2+71x+29)\,x^3$\\
79 & $(x+61)(x^2+70x+51)\,x^3$\\
83 & $(x+67)(x^2+9x+79)\,x^3$\\
89 & $(x^2+14x+78)(x^2+9x+25)\,x^3$\\
97 & $(x+69)(x^3+50x^2+8x+33)\,x^4$\\
101 & $(x+27)(x^3+60x^2+22x+59)\,x^4$\\
103 & $(x+90)(x^3+35x^2+61x+81)\,x^4$\\
107 & $(x^4+47x^3+31x^2+23x+36)\,x^4$\\
109 & $(x^5+61x^4+20x^3+35x^2+42x+40)\,x^4$\\
113 & $(x+41)(x+50)(x^3+81x^2+56)\,x^4$\\
\bottomrule
\end{tabular}
\end{table}
For $\ell=13,17,19,23$, the eigenvalue $-c$ of the factor $x+c$ is $\kappa^{(0)}_\ell$: $-7\equiv6$, $-11\equiv6$, $-9\equiv10$ and
$-18\equiv5$. For $\ell\ge59$ some unit eigenvalues are irreducible over $\Fl$, so for rank $\ge2$ the analogue of
``$L(b+2)\equiv\kappa L(b)$'' is a linear recurrence of order $R_\ell$.

\paragraph{$\ell=859$} \st{Computed} with method (M4), and confirmed independently by the $\ell^2R$ method. Here $d=71$ and $R=36$.
The results:
\begin{itemize}
\item $\det(\Xb|_\cV)=318\bmod859$, so $\rank\Xb|_\cV=36$;
\item $\rank\Db=36$ on both $\Sb$ and $\Mb$;
\item $\rank\Xb^k=36$ for $k\le72$;
\item $d_{859}=0$ by the exact test $\rank[\Xb^d\mid\Db_M]=\rank\Xb^d$;
\item the Krylov rank of $\Lb(1)$ is $36$, and $\charpoly(\Xb|_\cV)$ is squarefree.
\end{itemize}
So at the first prime where $6.6.a.a$ is known to be non-ordinary, Theorem~\ref{thm:B}, $d_\ell=0$ and rank attainment all hold.

% =====================================================================
\section{Open problems}\label{sec:open}
% =====================================================================

\begin{enumerate}[label=(\arabic*)]
\item \textbf{Sharpness for all $\ell\ge29$.} \st{Conjecture.} For every prime $\ell\ge29$ and every $m\ge1$, $b_\ell(m)=2m-1$. This would follow from
  (a) the slope-$1$ count: $X$ has exactly $\lfloor d/2\rfloor$ slopes equal to $1$, and $\Theta$ is invertible (gap G4); together with
  (b) $\bar u_0\ne0$, a nonvanishing statement for one explicit vector (gap G5). A possible route to (a) uses Ramsey and
  Coleman~\cite{FKO2012,Coleman1996}: the ordinary part of a $\theta$-twisted space of weight $\ell-5$ and level $6$.
  We have not worked this out.
\item \textbf{Equality $b_\ell(m)=m$ for $\ell\in\{13,17,19,23\}$.} Theorem~\ref{thm:A} gives $b_\ell(m)\le m$; equality is computed for
  $m\le16$. A proof of the lower bound requires the nonvanishing of the slope-$2$ component.
\item \textbf{Rank at higher precision (G6).} Show that $r_\ell(m)$ equals the Krylov rank of $\Lb(1)$ for every $m$, or find a
  counterexample.
\item \textbf{The exceptional set.} Is $\cE=\{\ell:\Lb(1)\text{ is not cyclic for }T_{\ell^2}\bmod\ell\}$ infinite? What is
  the second element after $607$? Is non-cyclicity always caused by a repeated $T_{\ell^2}$-eigenvalue in practice?
\item \textbf{Coloured partitions.} Extend Theorem~\ref{thm:B} and Corollary~\ref{cor:dl} to the BLMT invariant $d_\ell(r)$ for
  $\ell\ge r+5$~\cite{BLMT2014}.
\item \textbf{Formalization of Theorem~\ref{thm:B}.} This would require half-integral-weight forms and Yang's isomorphism (or a substitute)
  in Mathlib. A more modest target is the formalization of the finite computations at a fixed $\ell$.
\item \textbf{Explicit $\kappa_\ell$.} Express $\kappa_\ell\in\Zl^\times$ in terms of the unit-root $U_{\ell^2}$-eigenvalue on overconvergent
  half-integral weight forms (Ramsey's appendix to~\cite{FKO2012}). Is there a closed formula for $\kappa_\ell\bmod\ell^2$?
\end{enumerate}

\bibliographystyle{amsplain}
\bibliography{refs}

% =====================================================================
% TODO list (keep in sync with the report)
%  1. Conventions: check Definition 1.1 against BW Sec. 1 / Thm 1.4 / Remark 1.
%  2. Novelty of l = 17, 19, 23 (Atkin 1968; specialists).
%  3. Unchecked citations: Katz 2.6.2, 3.11.2, 3.11.3, App. 2; Coleman p. 221; Yang published version;
%     Deligne--Serre lemma number; theta-cycle references (G2); all page ranges marked TODO in refs.bib.
%  4. Garvan--Shomanov, El-Guindy--Ghazy, Webb thesis: complete the bib entries.
%  5. Lean: pin commit; run lean4checker; package UB/DB/HQ for all l >= 13.
%  6. G6 status: re-read docs/BL_SHARPNESS_PROOF.md (a gap-closing pass is running in parallel).
%  7. Truncation justification for Computational result 6.2 (theta, PR = 3, N = 6).
%  8. l = 607: BW for all m or m = 1; one vs two Galois orbits.
%  9. Slope counts for 181 <= l <= 229 at higher precision.
% 10. Rename dim S_{l-1} to avoid the clash with the BW invariant d_l.
% 11. Authorship, acknowledgements, AI-assistance disclosure.
% =====================================================================
\end{document}
